% Options for packages loaded elsewhere
\PassOptionsToPackage{unicode}{hyperref}
\PassOptionsToPackage{hyphens}{url}
%
\documentclass[
]{article}
\usepackage{amsmath,amssymb}
\usepackage{lmodern}
\usepackage{iftex}
\ifPDFTeX
  \usepackage[T1]{fontenc}
  \usepackage[utf8]{inputenc}
  \usepackage{textcomp} % provide euro and other symbols
\else % if luatex or xetex
  \usepackage{unicode-math}
  \defaultfontfeatures{Scale=MatchLowercase}
  \defaultfontfeatures[\rmfamily]{Ligatures=TeX,Scale=1}
\fi
% Use upquote if available, for straight quotes in verbatim environments
\IfFileExists{upquote.sty}{\usepackage{upquote}}{}
\IfFileExists{microtype.sty}{% use microtype if available
  \usepackage[]{microtype}
  \UseMicrotypeSet[protrusion]{basicmath} % disable protrusion for tt fonts
}{}
\makeatletter
\@ifundefined{KOMAClassName}{% if non-KOMA class
  \IfFileExists{parskip.sty}{%
    \usepackage{parskip}
  }{% else
    \setlength{\parindent}{0pt}
    \setlength{\parskip}{6pt plus 2pt minus 1pt}}
}{% if KOMA class
  \KOMAoptions{parskip=half}}
\makeatother
\usepackage{xcolor}
\usepackage{longtable,booktabs,array}
\usepackage{multirow}
\usepackage{calc} % for calculating minipage widths
% Correct order of tables after \paragraph or \subparagraph
\usepackage{etoolbox}
\makeatletter
\patchcmd\longtable{\par}{\if@noskipsec\mbox{}\fi\par}{}{}
\makeatother
% Allow footnotes in longtable head/foot
\IfFileExists{footnotehyper.sty}{\usepackage{footnotehyper}}{\usepackage{footnote}}
\makesavenoteenv{longtable}
\usepackage{graphicx}
\makeatletter
\def\maxwidth{\ifdim\Gin@nat@width>\linewidth\linewidth\else\Gin@nat@width\fi}
\def\maxheight{\ifdim\Gin@nat@height>\textheight\textheight\else\Gin@nat@height\fi}
\makeatother
% Scale images if necessary, so that they will not overflow the page
% margins by default, and it is still possible to overwrite the defaults
% using explicit options in \includegraphics[width, height, ...]{}
\setkeys{Gin}{width=\maxwidth,height=\maxheight,keepaspectratio}
% Set default figure placement to htbp
\makeatletter
\def\fps@figure{htbp}
\makeatother
\usepackage[normalem]{ulem}
\setlength{\emergencystretch}{3em} % prevent overfull lines
\providecommand{\tightlist}{%
  \setlength{\itemsep}{0pt}\setlength{\parskip}{0pt}}
\setcounter{secnumdepth}{-\maxdimen} % remove section numbering
\ifLuaTeX
  \usepackage{selnolig}  % disable illegal ligatures
\fi
\IfFileExists{bookmark.sty}{\usepackage{bookmark}}{\usepackage{hyperref}}
\IfFileExists{xurl.sty}{\usepackage{xurl}}{} % add URL line breaks if available
\urlstyle{same} % disable monospaced font for URLs
\hypersetup{
  hidelinks,
  pdfcreator={LaTeX via pandoc}}

\author{}
\date{}

\begin{document}

\begin{longtable}[]{@{}
  >{\raggedright\arraybackslash}p{(\columnwidth - 4\tabcolsep) * \real{0.3333}}
  >{\raggedright\arraybackslash}p{(\columnwidth - 4\tabcolsep) * \real{0.3333}}
  >{\raggedright\arraybackslash}p{(\columnwidth - 4\tabcolsep) * \real{0.3333}}@{}}
\toprule()
\begin{minipage}[b]{\linewidth}\raggedright
  \includegraphics[width=1.26528in,height=0.41528in]{vertopal_76ea1cc15b3d4460874f5e09e8e3ea79/media/image1.png}
\end{minipage} & \begin{minipage}[b]{\linewidth}\raggedright
\end{minipage} & \begin{minipage}[b]{\linewidth}\raggedright
\begin{quote}
2026-05-19
\end{quote}
\end{minipage} \\
\midrule()
\endhead
\multirow{2}{*}{\includegraphics[width=0.37083in,height=0.39861in]{vertopal_76ea1cc15b3d4460874f5e09e8e3ea79/media/image2.png}}
& \begin{minipage}[t]{\linewidth}\raggedright
\begin{quote}
\textbf{QuantaAlpha: An Evolutionary Framework for}
\end{quote}
\end{minipage} & \multirow{2}{*}{} \\
& \begin{minipage}[t]{\linewidth}\raggedright
\begin{quote}
\textbf{LLM-Driven Alpha Mining}
\end{quote}
\end{minipage} \\
\bottomrule()
\end{longtable}

\begin{quote}
\textbf{Jun Han1*, Shuo Zhang2*, Wei Li1*, Yifan Dong1, Tu Hu2, Yumo
Zhu1,}\\
\textbf{Xiaomin Yu2}, \textbf{Xin Guo1}, \textbf{Zhaowei Liu1},
\textbf{Kunyi Wang2}, \textbf{Jingping Liu3}, \textbf{Tianyi Jiang4},
\textbf{Ruichuan An4}, \textbf{Sen Hu2,4}, \textbf{Zhi Yang1†},
\textbf{Ronghao Chen2,4†}, \textbf{Huacan Wang2†}
\end{quote}

1SUFE, 2QuantaAlpha, 3SYSU, 4PKU

\textbf{*These authors contributed equally to this work.}

\begin{quote}
\textbf{†Correspondence:},,
\end{quote}

\begin{longtable}[]{@{}
  >{\raggedright\arraybackslash}p{(\columnwidth - 4\tabcolsep) * \real{0.3333}}
  >{\raggedright\arraybackslash}p{(\columnwidth - 4\tabcolsep) * \real{0.3333}}
  >{\raggedright\arraybackslash}p{(\columnwidth - 4\tabcolsep) * \real{0.3333}}@{}}
\toprule()
\begin{minipage}[b]{\linewidth}\raggedright
\end{minipage} &
\multicolumn{2}{>{\raggedright\arraybackslash}p{(\columnwidth - 4\tabcolsep) * \real{0.6667} + 2\tabcolsep}@{}}{%
\begin{minipage}[b]{\linewidth}\raggedright
\end{minipage}} \\
\midrule()
\endhead
&
\multicolumn{2}{>{\raggedright\arraybackslash}p{(\columnwidth - 4\tabcolsep) * \real{0.6667} + 2\tabcolsep}@{}}{%
\begin{minipage}[t]{\linewidth}\raggedright
\textbf{Abstract}

\begin{quote}
Financial markets are noisy and non-stationary, making alpha mining
highly sensitive to backtest noise and regime shifts. While recent
agentic frameworks improve automation, they often lack controllable
multi-round search and reliable reuse of validated experience. To
address these challenges, we propose \textbf{QuantaAlpha}, an
evolutionary alpha mining framework that treats each end-to-end mining
run as a trajectory and improves factors via trajectory-level mutation
and crossover. QuantaAlpha localizes suboptimal steps for targeted
revision and recombines complementary high-reward segments to reuse
effective patterns, enabling structured exploration and refinement
across iterations. During factor generation, it enforces semantic
consistency across hypothesis, factor expression, and executable code,
and constrains the complexity and redundancy of the generated factor to
mitigate crowding. Extensive experiments on CSI 300 show consistent
gains over strong baselines and prior agentic systems. Using GPT-5.2,
QuantaAlpha achieves an IC of 0.0472 with ARR of 4.68\% and MDD of
11.8\%. Moreover, factors mined on CSI 300 transfer effectively to CSI
500 and the S\&P 500, delivering about 40.28\% and 19.1\% cumulative
excess return over four years, respectively, which indicates strong
robustness under market distribution shifts.
\end{quote}
\end{minipage}} \\
&
\includegraphics[width=3.80694in,height=2.73056in]{vertopal_76ea1cc15b3d4460874f5e09e8e3ea79/media/image3.png}
& \begin{minipage}[t]{\linewidth}\raggedright
\begin{quote}
\includegraphics[width=2.15694in,height=2.73056in]{vertopal_76ea1cc15b3d4460874f5e09e8e3ea79/media/image4.png}
\end{quote}
\end{minipage} \\
\bottomrule()
\end{longtable}

\begin{quote}
Figure 1: Overview and empirical performance of QuantaAlpha. Left:
QuantaAlpha improves alpha discovery through trajectory-level
self-evolution. Right: Cumulative excess returns of different approaches
on CSI 500 and S\&P 500.
\end{quote}

1

\begin{longtable}[]{@{}
  >{\raggedright\arraybackslash}p{(\columnwidth - 2\tabcolsep) * \real{0.5000}}
  >{\raggedright\arraybackslash}p{(\columnwidth - 2\tabcolsep) * \real{0.5000}}@{}}
\toprule()
\begin{minipage}[b]{\linewidth}\raggedright
\begin{quote}
\textbf{1}
\end{quote}
\end{minipage} & \begin{minipage}[b]{\linewidth}\raggedright
\begin{quote}
\textbf{Introduction}
\end{quote}
\end{minipage} \\
\midrule()
\endhead
\bottomrule()
\end{longtable}

\begin{quote}
Financial markets are high-dimensional, non-stationary stochastic
systems, where returns exhibit heavy tails (Fama, 1965), time-varying
volatility (Engle, 1982), and cross-sectional dependence (Pesaran,
2021). Quantitative investment therefore relies on alpha mining to
extract predictive signals from noisy data. Recently, large language
models (LLMs) (Lopez-Lira and Tang, 2023) and LLM-based agent frameworks
(Xiao et al., 2024; Zhang et al., 2024) have been introduced into factor
research. By leveraging reasoning and code generation capabilities,
these systems can automate factor construction and iteratively refine
candidates through backtesting feedback.

Representative agent frameworks for alpha mining mimic human
quantitative researchers by iteratively performing hypothesis
generation, factor construction, and backtesting-based refinement. Under
this paradigm, AlphaAgent (Tang et al., 2025b) imposes explicit
regularization to mitigate crowding and alpha decay, while RD-Agent (Li
et al., 2025d) enables full-stack automation through coordinated
research and development agents for joint factor--model optimization.
These frameworks reduce trial-and-error costs while preserving
interpretability, making large-scale exploration practical.
\end{quote}

However, under real-market non-stationarity and low signal-to-noise
ratios, existing systems still face three limitations. \emph{(i) Fragile
controllability:} Iterative refinement is driven by noisy backtest
feedback, which can induce semantic drift and steer updates toward
spurious correlations, gradually deviating from the intended economic
mechanism. \emph{(ii) Limited trustworthiness:} Many methods rely on
stochastic re-generation conditioned on transient contexts, without
explicitly inheriting validated rationales across iterations. As a
result, they lack a traceable lineage and the produced factors are
harder to audit and trust. \emph{(iii) Constrained exploration:} Search
often over-exploits local neighborhoods around initial seeds, leading to
redundancy and factor crowding, while providing limited systematic
coverage of the broader hypothesis space and weakening long-horizon
discovery.

\begin{quote}
To address these challenges, we propose QuantaAlpha, an evolutionary
alpha mining framework that improves factor quality through
trajectory-level self-evolution (Figure 1). Inspired by (Lin et al.,
2025), we view each end-to-end mining run as a mining trajectory and
explicitly evolve trajectories rather than relying on unconstrained
re-generation from noisy feedback. First, to mitigate local
optimization, we introduce a Diversified Planning Initialization that
generates multiple complementary research directions, yielding broad
initial coverage in the hypothesis space. Second, to support
controllable refinement and reliable reuse of validated experience,
QuantaAlpha performs self-evolution via trajectory-level mutation and
crossover. Mutation performs targeted revision by localizing the
failure-causing step via self-reflection and rewriting only the
corresponding segment, while keeping the rest of the trajectory intact.
Crossover recombines complementary segments from high-reward parents to
reuse validated hypothesis structures, construction patterns, and repair
behaviors. We further apply generation gates to enforce semantic
consistency and constrain complexity and redundancy, preventing drift
and crowding. Together, these mechanisms expand exploration while
stabilizing refinement under real-market noise and non-stationarity.
Extensive experiments on CSI 300 and cross-market transfer to CSI 500
and S\&P 500 demonstrate consistent improvements in predictive power and
strategy performance over strong baselines and prior agentic systems.

\textbf{2 Related Work}
\end{quote}

\textbf{Agents in Finance.} Recent advances in financial LLMs (Liu et
al., 2023, 2025b) and evaluation bench-marks (Guo et al., 2025, 2026;
Tang et al., 2025a) have substantially expanded the scope of automated
financial reasoning. Building on these foundations, a growing line of
work explores how agentic systems (Li et al., 2025a; Yang et al., 2026)
can move beyond financial understanding toward decision-oriented
research workflows, including factor discovery and trading analysis. In
this context, the pursuit of alpha factors has evolved from manual
engineering and heuristic search toward LLM-driven closed-loop
discov-ery (Chen et al., 2025; Li et al., 2024b; Shi et al., 2025b; Wang
et al., 2025). Beyond search optimization, recent efforts emphasize
workflow systematization. RD-Agent (Li et al., 2025d) proposes a
multi-agent framework that decouples the pipeline into research and
development stages, enabling data-centric joint optimization of factors
and models. To address market non-stationarity, AlphaForge (Shi et al.,
2025a)

2

\begin{quote}
focuses on the dynamic combination of mined factors, while Alphafin (Li
et al., 2024a) and AlphaEval (Ding et al., 2025) establish standardized,
task-oriented protocols for reproducible evaluation. Despite these
advances, trading constraints (e.g., turnover, complexity) remain
largely post-hoc filters rather than intrinsic objectives, leading to
suboptimal generalization and interpretability in live trading.

\textbf{Self-Evolving Agents.} Self-evolving agents (Fang et al., 2025;
Lin et al., 2025; Zhai et al., 2025; Zhang et al., 2026) represent a
paradigm shift from static instruction-following to autonomous learning
through interacting with environment. AlphaEvolve (Novikov et al., 2025)
demonstrates a coding-centric evolution where the agent employs
evolutionary resampling to autonomously generate algorithms for
scientific discovery. Building on this, CSE (Hu et al., 2026) introduces
controllable self-evolution, facilitating a critical transition from
stochastic generation to feedback-driven evolution, while CogAlpha (Liu
et al., 2025a) applies code-level evolutionary search to alpha mining.
The financial domain adapts this self-evolutionary framework to handle
non-stationary and high-dimensional data. Research has converged on
multi-agent coordination and structured feedback to guide evolution (Li
et al., 2025b). TradingAgents (Xiao et al., 2024) and FactorMAD (Duan et
al., 2025) utilize institutional-style debates to refine trading
hypotheses, while QuantAgents (Li et al., 2025c) and ATLAS (Papadakis et
al., 2025) incorporate simulated trading performance as a reward signal
for dynamic prompt optimization. To ensure consistency over long
horizons, FinMem (Yu et al., 2025) and FinCon (Yu et al., 2024) leverage
hierarchical memory and conceptual reinforcement to retain and refine
high-level trading experience. Despite these advances, applying
self-evolving agents to finance is hindered by low signal-to-noise
ratios and delayed market feedback. To mitigate this, QuantaAlpha
employs stable optimization units with mutation and crossover operators,
ensuring a robust, traceable evolution path through structured
archiving.

\textbf{3 Problem Setup}
\end{quote}

\textbf{Alpha Mining.} The alpha mining task aims to learn an alpha
factor \emph{f} from a market feature tensor \textbf{X}
\emph{$\in$}R\emph{N×T×D}for a stock universe \emph{S} = \emph{\{s}1\emph{,
. . . , sN\}} and a time horizon \emph{T} = \emph{\{t}1\emph{, . . . ,
tT \}}. At each time \emph{t $\in$T} , the factor produces a signal from the
feature slice \textbf{X}\emph{t $\in$}R\emph{N×D}that is used to predict the
cross-sectional return \textbf{y}\emph{t}+1 \emph{$\in$}R\emph{N}. In
notation, an alpha can be written as \emph{f}(\textbf{X}\emph{t})
\emph{→}\textbf{y}\emph{t}+1, where \textbf{y}\emph{t}+1 denotes the
realized returns at time \emph{t} + 1. Accordingly, we formulate alpha
mining as the following optimization:

\begin{quote}
\emph{f$*$}= arg max \emph{f$\in$FL} \emph{f}(\textbf{X})\emph{,} \textbf{y}
\emph{$-$$\lambda$R}(\emph{f})\emph{,} (1)

where \emph{F} denotes the space of all possible factor expressions,
\textbf{y} is the ground-truth return target (e.g., next-day returns),
\emph{L}(\emph{·}) measures predictive effectiveness, \emph{R}(\emph{·})
is a regularization term that encourages simplicity and novelty of the
expression, and \emph{$\lambda$} balances utility and regularization.

\textbf{Alpha Mining Trajectory.} Distinct from purely data-driven
pipelines, we introduce market hypotheses \emph{h $\in$H} to guide LLM-based
factor construction. Each alpha mining run follows an end-to-end
workflow from hypothesis generation to factor generation and backtesting
evaluation (see Section 4.1). We represent each run by a mining
trajectory, defined as an ordered sequence \emph{$\tau$} = (\emph{s}0\emph{,
a}0\emph{, s}1\emph{, a}1\emph{, . . . , sn}), where \emph{s}0 denotes
the initial mining context (e.g., market context and optional
user-provided seeds), \emph{ai} is the action taken at step \emph{i} by
the multi-agent system, and \emph{sn} is the terminal state containing
the evaluated result of this run. The quality of a trajectory is
measured by its terminal reward:
\end{quote}

\begin{longtable}[]{@{}
  >{\raggedright\arraybackslash}p{(\columnwidth - 6\tabcolsep) * \real{0.2500}}
  >{\raggedright\arraybackslash}p{(\columnwidth - 6\tabcolsep) * \real{0.2500}}
  >{\raggedright\arraybackslash}p{(\columnwidth - 6\tabcolsep) * \real{0.2500}}
  >{\raggedright\arraybackslash}p{(\columnwidth - 6\tabcolsep) * \real{0.2500}}@{}}
\toprule()
\begin{minipage}[b]{\linewidth}\raggedright
\emph{R}(\emph{$\tau$}) = \emph{L}
\end{minipage} & \begin{minipage}[b]{\linewidth}\raggedright
\emph{f$\tau$}(\textbf{X})\emph{,} \textbf{y}
\end{minipage} & \begin{minipage}[b]{\linewidth}\raggedright
\begin{quote}
\emph{$-$$\lambda$R}(\emph{f$\tau$})\emph{,}
\end{quote}
\end{minipage} & \begin{minipage}[b]{\linewidth}\raggedright
(2)
\end{minipage} \\
\midrule()
\endhead
\bottomrule()
\end{longtable}

\begin{quote}
where \emph{f$\tau$} denotes the factor produced by trajectory \emph{$\tau$}.

\textbf{Objective.} Our objective is to learn a trajectory generation
policy \emph{$\pi$} for the multi-agent system that maximizes the expected
terminal reward of the generated mining trajectory:
\end{quote}

\begin{longtable}[]{@{}
  >{\raggedright\arraybackslash}p{(\columnwidth - 8\tabcolsep) * \real{0.2000}}
  >{\raggedright\arraybackslash}p{(\columnwidth - 8\tabcolsep) * \real{0.2000}}
  >{\raggedright\arraybackslash}p{(\columnwidth - 8\tabcolsep) * \real{0.2000}}
  >{\raggedright\arraybackslash}p{(\columnwidth - 8\tabcolsep) * \real{0.2000}}
  >{\raggedright\arraybackslash}p{(\columnwidth - 8\tabcolsep) * \real{0.2000}}@{}}
\toprule()
\begin{minipage}[b]{\linewidth}\raggedright
\begin{quote}
\emph{$\pi$$*$}= arg max\\
\emph{$\pi$}
\end{quote}\strut
\end{minipage} & \begin{minipage}[b]{\linewidth}\raggedright
E\emph{$\tau$$\sim$$\pi$}
\end{minipage} & \begin{minipage}[b]{\linewidth}\raggedright
\emph{R}(\emph{$\tau$})
\end{minipage} & \begin{minipage}[b]{\linewidth}\raggedright
\begin{quote}
\emph{,}
\end{quote}
\end{minipage} & \begin{minipage}[b]{\linewidth}\raggedright
(3)
\end{minipage} \\
\midrule()
\endhead
\bottomrule()
\end{longtable}

\begin{quote}
where \emph{$\tau$ $\sim$$\pi$} denotes the trajectory induced by repeatedly applying
\emph{$\pi$} from the initial state \emph{s}0.
\end{quote}

3

\begin{longtable}[]{@{}
  >{\raggedright\arraybackslash}p{(\columnwidth - 2\tabcolsep) * \real{0.5000}}
  >{\raggedright\arraybackslash}p{(\columnwidth - 2\tabcolsep) * \real{0.5000}}@{}}
\toprule()
\begin{minipage}[b]{\linewidth}\raggedright
\begin{quote}
\textbf{4}
\end{quote}
\end{minipage} & \begin{minipage}[b]{\linewidth}\raggedright
\begin{quote}
\textbf{Method}
\end{quote}
\end{minipage} \\
\midrule()
\endhead
\bottomrule()
\end{longtable}

\begin{quote}
As illustrated in Figure 2, our approach treats alpha mining as an
agentic workflow rather than one-shot factor construction. Given market
context and user-provided seeds, a multi-agent system produces a mining
trajectory including hypothesis generation, factor construction,
implementation, and backtesting evaluation. We first describe the
trajectory generation workflow (Section 4.1), and then introduce a
self-evolution scheme that iteratively improves mining quality (Section
4.2).
\end{quote}

\includegraphics[width=5.66944in,height=3.02361in]{vertopal_76ea1cc15b3d4460874f5e09e8e3ea79/media/image5.png}

\begin{quote}
Figure 2: Overview of the QuantaAlpha framework, which consists of four
core components: (A) Diversified Planning Initialization for candidate
hypotheses, (B) Factor Realization that instantiates hypotheses into
executable factors with constraint gating, (C) Self-Evolution with
mutation and crossover over evaluated trajectories, and (D) A Final
Factor Pool that consolidates validated factors.

\textbf{4.1} \textbf{Alpha Mining Process}

A standalone alpha mining process progresses from hypothesis generation
to factor creation and backtesting-based evaluation. We translate this
process into collaboration among multiple agents, each responsible for a
specific task at different stages of the workflow.

\textbf{4.1.1 Hypothesis Generation}
\end{quote}

We use the idea agent \emph{Ai} to generate market hypotheses via
structured knowledge integration. Conditioned on (i) observations from
current market conditions or previous mining trajectories, (ii) domain
priors from financial theories and empirical evidence, and (iii)
parametric specifications (e.g., ``10-day high/low''), \emph{Ai}
produces actionable hypotheses that describe candidate market mechanisms
and serve as the input to subsequent factor generation.

\begin{quote}
\textbf{4.1.2 Controllable Factor Construction}

To effectively instantiate hypothesis-driven factors, we need a robust
mechanism to align implementations with domain hypotheses. However,
direct code generation is brittle, often suffering from syntax errors,
dependency mismatches, and semantic drift in complex implementations,
which undermines fidelity to the hypothesis. To address these
limitations, We introduce an intermediate symbolic representation based
on a standardized operator library \emph{O} and an Abstract Syntax Tree
(AST). This abstraction bridges high-level market intent and low-level
implementation, enabling controllable construction, structural
validation, and reliable compilation.

\textbf{Factor Generation.} Given a hypothesis \emph{h $\in$H} and raw
features \emph{X}, the factor agent \emph{Af} generates a symbolic
expression \emph{f $\in$F} over the operator library and parses it into an
AST as the intermediate representation. Concretely, it first maps
\emph{h} to a structured semantic description \emph{d} that formalizes
the
\end{quote}

4

\begin{quote}
intended mechanism into discrete operators, while concurrently
instantiating operational parameters---such as look-back windows and
thresholds---from either prescribed parameters in \emph{h} or heuristic
defaults anchored in domain expertise. Conditioned on \emph{d}, the
agent then assembles an operator composition into a symbolic expression
\emph{f} and parses it into an AST, denoted as \emph{T} (\emph{f}). In
\emph{T}(\emph{f}), leaf nodes bind to raw feature fields (e.g., \$high,
\$volume) and internal nodes correspond to operator instances from
\emph{O} (e.g., TS\_MIN (\emph{·}), SMA (\emph{·}), RANK (\emph{·})),
thereby rendering the computational dependencies and data flow fully
transparent. Finally, the agent translates the validated symbolic form
into executable code \emph{c} via compilation; when compilation fails
due to implementation issues, an LLM-based repair step is triggered to
regenerate a consistent implementation while preserving the semantics of
\emph{f}. This design ensures that code generation remains anchored to
an explicit symbolic specification rather than unconstrained free-form
generation.

\textbf{Consistency Verification.} We enforce semantic consistency
across representations to prevent drift during generation. Specifically,
we apply an LLM-based verifier to assess (i) alignment among the
hypothesis \emph{h}, semantic description \emph{d}, and symbolic
expression \emph{f}, and (ii) faithfulness between the symbolic
definition \emph{f} and the generated code \emph{c}. The verifier
returns a consistency decision under fixed criteria. If it fails, we
rewrite the inconsistent component(s) by regenerating \emph{d} and
\emph{f} or repairing \emph{c} until the check passes or a maximum retry
budget is reached.

\textbf{Complexity and Redundancy Control.} We further regularize factor
generation with explicit structural constraints to promote parsimony and
novelty. For complexity, we compute
\end{quote}

\begin{longtable}[]{@{}
  >{\raggedright\arraybackslash}p{(\columnwidth - 6\tabcolsep) * \real{0.2500}}
  >{\raggedright\arraybackslash}p{(\columnwidth - 6\tabcolsep) * \real{0.2500}}
  >{\raggedright\arraybackslash}p{(\columnwidth - 6\tabcolsep) * \real{0.2500}}
  >{\raggedright\arraybackslash}p{(\columnwidth - 6\tabcolsep) * \real{0.2500}}@{}}
\toprule()
\begin{minipage}[b]{\linewidth}\raggedright
\begin{quote}
\emph{C}(\emph{f}) = \emph{$\alpha$}1 \emph{· SL}(\emph{f}) + \emph{$\alpha$}2 \emph{·
PC}(\emph{f}) + \emph{$\alpha$}3 \emph{·} log
\end{quote}
\end{minipage} & \begin{minipage}[b]{\linewidth}\raggedright
1 + \emph{\textbar Ff\textbar{}}
\end{minipage} & \begin{minipage}[b]{\linewidth}\raggedright
\begin{quote}
\emph{,}
\end{quote}
\end{minipage} & \begin{minipage}[b]{\linewidth}\raggedright
(4)
\end{minipage} \\
\midrule()
\endhead
\bottomrule()
\end{longtable}

\begin{quote}
where \emph{SL}(\emph{f}) measures symbolic length, \emph{PC} (\emph{f})
counts free parameters (e.g., window sizes), and \emph{Ff} is the set of
features used by \emph{f}. For redundancy, we quantify structural
similarity via AST matching. Given two factors \emph{fi} and \emph{fj}
with ASTs \emph{T}(\emph{fi}) and \emph{T}(\emph{fj}), we define
\emph{s}(\emph{fi, fj}) as the size of the largest common subtree:

\emph{s}(\emph{fi, fj}) =
\emph{Si$\subseteq$T}(\emph{fi})\emph{,Sj$\subseteq$T}(\emph{fj})\emph{,Si$\sim$} max =\emph{Sj
\textbar V} (\emph{Si})\emph{\textbar.} (5)

Here, \emph{Si} and \emph{Sj} range over subtrees of \emph{T}(\emph{fi})
and \emph{T}(\emph{fj}), \emph{Si$\sim$}= \emph{Sj} denotes subtree
isomorphism, and \emph{\textbar V} (\emph{Si})\emph{\textbar{}} is the
number of nodes in \emph{Si}. For a candidate factor \emph{f}, we
compute its maximum similarity to an existing alpha zoo \emph{Z} as
\emph{S}(\emph{f}) = max\emph{$\phi$$\in$Z s}(\emph{f, $\phi$})\emph{.} We reject and
rewrite any factor that violates the prescribed complexity or redundancy
limits, pro-moting parsimonious generation while avoiding near-duplicate
structures. We further apply output-level correlation filtering on the
2021 validation panel to remove functionally equivalent factors: if two
factors exceed an absolute correlation threshold, the one with lower
RankIC is discarded. This two-stage design handles both structural and
functional redundancy.

\textbf{4.1.3 Factor Evaluation}

The evaluation agent assesses factors via a standardized backtesting
system. The protocol covers three aspects: predictive capability metrics
for forecasting effectiveness, return performance metrics for
profit-generating potential under a fixed execution setting, and risk
control metrics for stability and robustness across conditions. Beyond
factor scoring, the agent maintains an evaluation history that records
outcomes and summarizes systematic patterns among successful and failed
factors.

\textbf{4.2} \textbf{Evolutionary Alpha Mining}
\end{quote}

We improve alpha discovery via iterative optimization over mining
trajectories. Starting from user-provided seed factors, we generate a
diversified set of initial hypotheses. We then run the workflow in
Section 4.1 to turn each hypothesis into an evaluated mining trajectory,
forming an initial trajectory pool \emph{T}0 = \emph{\{$\tau$}0
\emph{j\}Ninit j}=1. We then apply an evolution process to obtain
progressively improved factors.

\begin{quote}
\textbf{4.2.1 Diversified Planning Initialization}
\end{quote}

Let \emph{Z}0 denote the user-provided seed factors (or seed ideas). The
initial seed pool is derived from the public Alpha158(20) subset,
grouped by low correlation to ensure diverse starting directions. We
employ

5

\begin{longtable}[]{@{}
  >{\raggedright\arraybackslash}p{(\columnwidth - 0\tabcolsep) * \real{1.0000}}@{}}
\toprule()
\begin{minipage}[b]{\linewidth}\raggedright
an initialization agent \emph{A}0 to propose a diversified set of
initial hypotheses \emph{H}0= \emph{\{h}0 agent is instructed to
maximize complementarity among hypotheses at both semantic and
structural levels, 1\emph{, . . . , h}0 \emph{Ninit\}}. The
\end{minipage} \\
\midrule()
\endhead
\bottomrule()
\end{longtable}

\begin{quote}
e.g., varying signal sources (price vs. volume), time scales (short-term
vs. long-term), and mechanism types (momentum vs. mean-reversion or
regime-conditioned triggers). For each \emph{h}0 mining workflow in
Section 4.1 and obtain a mining trajectory \emph{$\tau$}0 \emph{i}.
\emph{i$\in$H}0, we run the

This diversified initialization expands the effective search frontier by
providing broad coverage in the hypothesis space. As a result, the
subsequent evolutionary process can explore multiple promising regions
in parallel, reducing the risk of premature convergence to a suboptimal
local optimum.

\textbf{4.2.2 Self-Evolution}
\end{quote}

The goal of self-evolutionary updates is to guide the alpha mining
system to search for better trajectories based on existing ones.
Specifically, we apply trajectory-level operators to existing mining
trajectories to generate improved trajectories as demonstrations. These
demonstrations provide imitation learning priors that bias subsequent
trajectory generation toward effective decisions. In practice, we
instantiate the update with two evolutionary primitives, Mutation and
Crossover, which revise a single trajectory or recombine complementary
sub-trajectories, respectively. Each iteration \emph{i} yields a new
trajectory generation \emph{Ti}, and factor quality improves
progressively across iterations (see Figure 6 in Appendix for a case
study).

sub-optimal decision node that most significantly accounts for the low
terminal reward, denoted by an \textbf{Mutation.} Given a mining
trajectory \emph{$\tau$ $\in$Ti$-$}1, the agent first performs self-reflection to
diagnose

index \emph{k $\in$\{}0\emph{, . . . , n $-$}1\emph{\}}. We then rewrite only
the localized action (or a short local segment) while keeping the
remaining steps unchanged:

\begin{quote}
\emph{$\tau$child} = \emph{s}0\emph{, a}0\emph{, . . . , sk,}
Refine(\emph{ak})\emph{, s$'$k}+1\emph{, a$'$k}+1\emph{, . . . , s$'$n}
\emph{,} (6)
\end{quote}

where the prefix up to \emph{sk} is frozen, and the remaining steps are
regenerated by the agent conditioned on this fixed prefix to preserve
trajectory coherence. Depending on the localized fault, the rewrite may
update the hypothesis, the symbolic expression under the operator
library, or the compiled code, and can introduce mechanism-level changes
such as altering the time scale, or adding regime conditions. The
diagnostic step provides a directional search signal, allowing even
imperfect localization to move the child trajectory toward a different
region of the factor space.

\begin{quote}
\textbf{Crossover.} Crossover aims to exploit validated components by
recombining complementary sub-trajectories from high-quality parents. At
iteration \emph{i $-$}1, we select a parent subset \emph{Pi$-$}1
\emph{$\subseteq$Ti$-$}1 based on trajectory reward (Eq. 2). Given \emph{k} parent
trajectories \emph{$\tau$}(1)\emph{, . . . , $\tau$}(\emph{k})\emph{$\in$Pi$-$}1,
Crossover synthesizes a child trajectory by composing high-performing
segments from different parents:

\emph{$\tau$}child = Crossover\emph{$\tau$}(1)\emph{, . . . , $\tau$}(\emph{k})
\emph{.} (7)

Operationally, the primitive identifies trajectory segments that
consistently contribute to high cumulative rewards---such as hypothesis
templates, factor construction patterns, or strategic repair
actions---and merges them into a highly coherent sequence. This
recombination mimics the human practice of combining complementary
strengths from different solutions, while providing a more credible
lineage by explicitly inheriting decisions validated in previously
successful trajectories.

\textbf{5 Experiments}

\textbf{5.1} \textbf{Experimental Setup}
\end{quote}

\textbf{Datasets and Metrics.} We conduct experiments on the CSI 300
dataset, which covers 300 large-cap A-share stocks in the Chinese
market. We adopt a chronological split with training (Jan. 1, 2016 to
Dec. 31, 2020), validation (Jan. 1, 2021 to Dec. 31, 2021), and testing
(Jan. 1, 2022 to Dec. 26, 2025). We evaluate performance from two
complementary perspectives. Factor predictive power is measured by
Information Coefficient (IC), IC Information Ratio (ICIR), Rank IC, and
Rank ICIR. Strategy performance is measured by Annualized Return (ARR),
Information Ratio (IR), and Maximum Drawdown (MDD). Further details are
provided in Appendix A.1.

6

\begin{quote}
\textbf{Baselines.} We compare QuantaAlpha against four baseline
categories: traditional machine learning models, deep learning
time-series models, classical factor libraries, and LLM-based factor
research agents, including RD-Agent and AlphaAgent. More details are
provided in Appendix A.3. For all LLM-based methods, roughly 150
validated factors are submitted to the same downstream LightGBM model
for final evaluation, ensuring a fair factor-pool-level comparison.

\textbf{5.2} \textbf{Main Results}
\end{quote}

Table 1 reports the main results on the CSI 300 market over a four-year
evaluation period, and strategy-level results are computed from a factor
pool of approximately 150 validated factors synthesized by the same
downstream LightGBM model, rather than from a single best factor. For
transparency regarding the search budget, we report computational cost
and token consumption in Appendix A.4. QuantaAlpha achieves the
strongest overall performance across most factor predictive power and
strategy-level performance. Using GPT-5.2, QuantaAlpha achieves the best
IC at 0.0472, delivers an ARR of 4.68\%, and maintains a low MDD of
11.80\%. From an empirical perspective, these results suggest that the
mined factors capture non-trivial predictive signals beyond simple
backtest noise, although real-world deployment would require additional
transaction-cost modeling, risk control, and live trading validation. We
also observe broadly consistent gains across different base models,
suggesting that the improvements are not tied to a single LLM backbone.
The narrower gap between Claude-4.5-Sonnet and DeepSeek-V3.2 within
QuantaAlpha (0.004 IC) is attributed to backbone-specific adaptation:
Claude's longer feedback outputs reduce the number of prior trajectories
fitting within the fixed context budget, and its mutations tend to be
more local, whereas DeepSeek-V3.2 more frequently proposes larger
exploratory jumps.

Compared with RD-Agent, QuantaAlpha and AlphaAgent both incorporate
generation-stage regu-larization, and this is reflected in stronger
predictive power and strategy performance. Under GPT-5.2, QuantaAlpha
improves IC by 0.0186 and ARR by 1.10\% relative to RD-Agent while
reducing MDD by 4.96\%, suggesting that constraining
hypothesis-to-factor construction with explicit intermediate
represen-tations and verification gates effectively reduces semantic
drift and improves implementation reliability. Beyond these constraints,
QuantaAlpha further improves upon AlphaAgent by 0.0125 IC and 3.57\% ARR
while reducing MDD by 2.09\%, highlighting the added value of
trajectory-centric self-evolution: mutation broadens exploration via
mechanism-level variations, while crossover reuses validated trajectory
segments and repair strategies, jointly increasing the yield and
stability of high-quality factors under non-stationary market dynamics.

\begin{quote}
\textbf{5.3} \textbf{Ablation Study}
\end{quote}

\textbf{Ablation of Evolutionary Mining Components.} We conduct an
ablation study to quantify the contri-bution of three components in
evolutionary alpha mining: diversified planning initialization,
trajectory mutation, and trajectory crossover. The results are
summarized in Table 2. Removing \emph{planning} yields only marginal
changes in IC and Rank IC, but substantially degrades strategy-level
outcomes, with ARR decreasing by 0.72\% and MDD increasing by 1.62\%.
This indicates that diversified initialization primarily improves the
search frontier by providing broader and less correlated starting
hypotheses, which stabilizes subsequent evolution. In contrast, removing
\emph{mutation} causes the largest drop in predictive power, decreasing
IC by 0.0079 and Rank IC by 0.0079, and also leads to the largest
reduction in ARR by 1.26\%. This highlights mutation as the primary
driver of effective exploration and trajectory repair, enabling the
system to escape suboptimal regions and correct failure modes discovered
during mining. Meanwhile, removing \emph{crossover} leads to a smaller
but consistent degradation. This supports the role of crossover in
exploiting and inheriting complementary high-performing trajectory
segments, improving efficiency and stability by recombining validated
patterns from successful trajectories. These results suggest that
mutation contributes more to exploration in the current five-cycle
setting, whereas crossover mainly supports reuse and recombination of
validated trajectory segments. More broadly, QuantaAlpha's contribution
lies in evolving complete research trajectories rather than relying on
any single operator.

\begin{quote}
\textbf{Ablation of Consistency, Complexity, and Redundancy Controls.}
We ablate three controls during factor generation that gate rewriting,
including semantic consistency verification, complexity regulariza-tion,
and redundancy filtering. As shown in Figure 3, disabling any single
control consistently degrades
\end{quote}

7

\begin{longtable}[]{@{}
  >{\raggedright\arraybackslash}p{(\columnwidth - 0\tabcolsep) * \real{1.0000}}@{}}
\toprule()
\begin{minipage}[b]{\linewidth}\raggedright
Table 1: Performance comparison of different methods on
\emph{\textbf{CSI 300}}. We report both \emph{factor predictive power}
and \emph{strategy-level performance} metrics, where
\textbf{higher}\emph{↑}values indicate better performance for all
metrics, except for MDD where lower values are better. For each method
category, the best result in each column is highlighted in
\textbf{bold}, while the second-best result is\uline{underlined}.
\end{minipage} \\
\midrule()
\endhead
\bottomrule()
\end{longtable}

\begin{quote}
Across all methods, the overall best-performing method is highlighted
with a light blue background .
\end{quote}

\begin{longtable}[]{@{}
  >{\raggedright\arraybackslash}p{(\columnwidth - 8\tabcolsep) * \real{0.2000}}
  >{\raggedright\arraybackslash}p{(\columnwidth - 8\tabcolsep) * \real{0.2000}}
  >{\raggedright\arraybackslash}p{(\columnwidth - 8\tabcolsep) * \real{0.2000}}
  >{\raggedright\arraybackslash}p{(\columnwidth - 8\tabcolsep) * \real{0.2000}}
  >{\raggedright\arraybackslash}p{(\columnwidth - 8\tabcolsep) * \real{0.2000}}@{}}
\toprule()
\begin{minipage}[b]{\linewidth}\raggedright
\textbf{Methods}
\end{minipage} & \begin{minipage}[b]{\linewidth}\raggedright
\end{minipage} &
\multicolumn{2}{>{\raggedright\arraybackslash}p{(\columnwidth - 8\tabcolsep) * \real{0.4000} + 2\tabcolsep}}{%
\begin{minipage}[b]{\linewidth}\raggedright
\begin{quote}
Factor Predictive Power
\end{quote}
\end{minipage}} & \begin{minipage}[b]{\linewidth}\raggedright
\begin{quote}
Strategy Performance
\end{quote}
\end{minipage} \\
\midrule()
\endhead
& \textbf{IC} & \textbf{ICIR} &
\multicolumn{2}{>{\raggedright\arraybackslash}p{(\columnwidth - 8\tabcolsep) * \real{0.4000} + 2\tabcolsep}@{}}{%
\textbf{Rank IC Rank ICIR IR (SHR*) ARR (\%) MDD (\%)}\emph{↓}} \\
\bottomrule()
\end{longtable}

\begin{longtable}[]{@{}
  >{\raggedright\arraybackslash}p{(\columnwidth - 16\tabcolsep) * \real{0.1111}}
  >{\raggedright\arraybackslash}p{(\columnwidth - 16\tabcolsep) * \real{0.1111}}
  >{\raggedright\arraybackslash}p{(\columnwidth - 16\tabcolsep) * \real{0.1111}}
  >{\raggedright\arraybackslash}p{(\columnwidth - 16\tabcolsep) * \real{0.1111}}
  >{\raggedright\arraybackslash}p{(\columnwidth - 16\tabcolsep) * \real{0.1111}}
  >{\raggedright\arraybackslash}p{(\columnwidth - 16\tabcolsep) * \real{0.1111}}
  >{\raggedright\arraybackslash}p{(\columnwidth - 16\tabcolsep) * \real{0.1111}}
  >{\raggedright\arraybackslash}p{(\columnwidth - 16\tabcolsep) * \real{0.1111}}
  >{\raggedright\arraybackslash}p{(\columnwidth - 16\tabcolsep) * \real{0.1111}}@{}}
\toprule()
\begin{minipage}[b]{\linewidth}\raggedright
\end{minipage} & \begin{minipage}[b]{\linewidth}\raggedright
\end{minipage} &
\multicolumn{3}{>{\raggedright\arraybackslash}p{(\columnwidth - 16\tabcolsep) * \real{0.3333} + 4\tabcolsep}}{%
\begin{minipage}[b]{\linewidth}\raggedright
\begin{minipage}[b]{\linewidth}\raggedright
\begin{quote}
\emph{\textbf{Classical Factor Mining}}
\end{quote}
\end{minipage}
\end{minipage}} & \begin{minipage}[b]{\linewidth}\raggedright
\end{minipage} & \begin{minipage}[b]{\linewidth}\raggedright
\end{minipage} & \begin{minipage}[b]{\linewidth}\raggedright
\end{minipage} & \begin{minipage}[b]{\linewidth}\raggedright
\end{minipage} \\
\midrule()
\endhead
\multirow{6}{*}{\textbf{Machine Learning}} &
\begin{minipage}[t]{\linewidth}\raggedright
\begin{quote}
Linear
\end{quote}
\end{minipage} & 0.0155 & 0.1174 &
\begin{minipage}[t]{\linewidth}\raggedright
\begin{quote}
0.0368
\end{quote}
\end{minipage} & 0.2834 & -0.3078 & -2.67 & 18.97 \\
& \begin{minipage}[t]{\linewidth}\raggedright
\begin{quote}
XGBoost
\end{quote}
\end{minipage} & 0.0175 & 0.1336 &
\begin{minipage}[t]{\linewidth}\raggedright
\begin{quote}
0.0420
\end{quote}
\end{minipage} & 0.3417 & -0.5280 & -4.24 & 28.50 \\
& \begin{minipage}[t]{\linewidth}\raggedright
\begin{quote}
CatBoost
\end{quote}
\end{minipage} & 0.0162 & 0.1203 &
\begin{minipage}[t]{\linewidth}\raggedright
\begin{quote}
0.0405
\end{quote}
\end{minipage} & 0.3289 & -0.2807 & -2.30 & 21.35 \\
& \begin{minipage}[t]{\linewidth}\raggedright
\begin{quote}
LightGBM
\end{quote}
\end{minipage} & \uline{0.0247} & \uline{0.2055} &
\begin{minipage}[t]{\linewidth}\raggedright
\begin{quote}
\uline{0.0423}
\end{quote}
\end{minipage} & \uline{0.3726} & 0.0092 & 0.07 & 21.80 \\
& \begin{minipage}[t]{\linewidth}\raggedright
\begin{quote}
MLP
\end{quote}
\end{minipage} & \textbf{0.0321} & \textbf{0.2780} &
\begin{minipage}[t]{\linewidth}\raggedright
\begin{quote}
\textbf{0.0438}
\end{quote}
\end{minipage} & \textbf{0.4088} & \uline{0.1716} & \uline{1.46} &
\uline{18.15} \\
& \begin{minipage}[t]{\linewidth}\raggedright
\begin{quote}
DoubleEnsemble
\end{quote}
\end{minipage} & 0.0213 & 0.1670 &
\begin{minipage}[t]{\linewidth}\raggedright
\begin{quote}
0.0408
\end{quote}
\end{minipage} & 0.3372 & \textbf{0.2490} & \textbf{1.85} &
\textbf{15.00} \\
\multirow{4}{*}{\begin{minipage}[t]{\linewidth}\raggedright
\begin{quote}
\textbf{Deep Learning}
\end{quote}
\end{minipage}} & \begin{minipage}[t]{\linewidth}\raggedright
\begin{quote}
GRU
\end{quote}
\end{minipage} & 0.0321 & 0.2603 &
\begin{minipage}[t]{\linewidth}\raggedright
\begin{quote}
0.0442
\end{quote}
\end{minipage} & 0.3601 & 0.5302 & 3.61 & 15.01 \\
& \begin{minipage}[t]{\linewidth}\raggedright
\begin{quote}
Transformer
\end{quote}
\end{minipage} & \uline{0.0331} & \uline{0.2702} &
\begin{minipage}[t]{\linewidth}\raggedright
\begin{quote}
\uline{0.0451}
\end{quote}
\end{minipage} & \uline{0.3801} & 0.4502 & 5.21 & \uline{13.81} \\
& \begin{minipage}[t]{\linewidth}\raggedright
\begin{quote}
LSTM
\end{quote}
\end{minipage} & \uline{0.0331} & 0.2502 &
\begin{minipage}[t]{\linewidth}\raggedright
\begin{quote}
\uline{0.0451}
\end{quote}
\end{minipage} & 0.3503 & \uline{0.6802} & \uline{6.01} & 14.81 \\
& \begin{minipage}[t]{\linewidth}\raggedright
\begin{quote}
TRA
\end{quote}
\end{minipage} & \textbf{0.0421} & \textbf{0.3402} &
\begin{minipage}[t]{\linewidth}\raggedright
\begin{quote}
\textbf{0.0511}
\end{quote}
\end{minipage} & \textbf{0.4203} & \textbf{1.0502} & \textbf{6.81} &
\textbf{8.51} \\
\multirow{3}{*}{\begin{minipage}[t]{\linewidth}\raggedright
\begin{quote}
\textbf{Factor Libraries}
\end{quote}
\end{minipage}} & \begin{minipage}[t]{\linewidth}\raggedright
\begin{quote}
Alpha158(20)
\end{quote}
\end{minipage} & 0.0051 & 0.0329 &
\begin{minipage}[t]{\linewidth}\raggedright
\begin{quote}
0.0184
\end{quote}
\end{minipage} & 0.1177 & \uline{0.5044} & \textbf{4.63} & 22.19 \\
& \begin{minipage}[t]{\linewidth}\raggedright
\begin{quote}
Alpha158
\end{quote}
\end{minipage} & \textbf{0.0131} & \textbf{0.0817} &
\begin{minipage}[t]{\linewidth}\raggedright
\begin{quote}
\textbf{0.0334}
\end{quote}
\end{minipage} & \textbf{0.2119} & 0.4099 & 2.66 & \textbf{10.15} \\
& \begin{minipage}[t]{\linewidth}\raggedright
\begin{quote}
Alpha360
\end{quote}
\end{minipage} & \uline{0.0105} & \uline{0.0636} &
\begin{minipage}[t]{\linewidth}\raggedright
\begin{quote}
\uline{0.0306}
\end{quote}
\end{minipage} & \uline{0.1889} & \textbf{0.6009} & \uline{4.09} &
\uline{11.52} \\
\bottomrule()
\end{longtable}

\begin{longtable}[]{@{}
  >{\raggedright\arraybackslash}p{(\columnwidth - 16\tabcolsep) * \real{0.1111}}
  >{\raggedright\arraybackslash}p{(\columnwidth - 16\tabcolsep) * \real{0.1111}}
  >{\raggedright\arraybackslash}p{(\columnwidth - 16\tabcolsep) * \real{0.1111}}
  >{\raggedright\arraybackslash}p{(\columnwidth - 16\tabcolsep) * \real{0.1111}}
  >{\raggedright\arraybackslash}p{(\columnwidth - 16\tabcolsep) * \real{0.1111}}
  >{\raggedright\arraybackslash}p{(\columnwidth - 16\tabcolsep) * \real{0.1111}}
  >{\raggedright\arraybackslash}p{(\columnwidth - 16\tabcolsep) * \real{0.1111}}
  >{\raggedright\arraybackslash}p{(\columnwidth - 16\tabcolsep) * \real{0.1111}}
  >{\raggedright\arraybackslash}p{(\columnwidth - 16\tabcolsep) * \real{0.1111}}@{}}
\toprule()
\begin{minipage}[b]{\linewidth}\raggedright
\end{minipage} & \begin{minipage}[b]{\linewidth}\raggedright
\end{minipage} &
\multicolumn{4}{>{\raggedright\arraybackslash}p{(\columnwidth - 16\tabcolsep) * \real{0.4444} + 6\tabcolsep}}{%
\begin{minipage}[b]{\linewidth}\raggedright
\begin{quote}
\emph{\textbf{LLM-based Agentic Factor Mining}}
\end{quote}
\end{minipage}} & \begin{minipage}[b]{\linewidth}\raggedright
\end{minipage} & \begin{minipage}[b]{\linewidth}\raggedright
\end{minipage} & \begin{minipage}[b]{\linewidth}\raggedright
\end{minipage} \\
\midrule()
\endhead
\multirow{5}{*}{\begin{minipage}[t]{\linewidth}\raggedright
\begin{quote}
\textbf{RD-Agent}
\end{quote}
\end{minipage}} & \begin{minipage}[t]{\linewidth}\raggedright
\begin{quote}
Qwen3-235B
\end{quote}
\end{minipage} & 0.0267 & 0.1676 & 0.0194 & 0.1199 & -0.0818 & -0.62 &
15.04 \\
& \begin{minipage}[t]{\linewidth}\raggedright
\begin{quote}
Deepseek-V3.2
\end{quote}
\end{minipage} & 0.0245 & 0.1630 & 0.0192 & 0.1250 & -0.2123 & -1.42 &
19.17 \\
& \begin{minipage}[t]{\linewidth}\raggedright
\begin{quote}
Gemini-3-pro-preview
\end{quote}
\end{minipage} & \textbf{0.0301} & 0.1870 & \textbf{0.0282} & 0.1677 &
0.2595 & 1.89 & \uline{11.49} \\
& Claude-4.5-sonnet & 0.0280 & \textbf{0.2000} & 0.0242 & \uline{0.1708}
& \uline{0.3568} & \uline{2.36} & \textbf{10.81} \\
& \begin{minipage}[t]{\linewidth}\raggedright
\begin{quote}
GPT-5.2
\end{quote}
\end{minipage} & \uline{0.0286} & \uline{0.1995} & \uline{0.0250} &
\textbf{0.1739} & \textbf{0.5321} & \textbf{3.58} & 16.76 \\
\multirow{5}{*}{\begin{minipage}[t]{\linewidth}\raggedright
\begin{quote}
\textbf{AlphaAgent}
\end{quote}
\end{minipage}} & \begin{minipage}[t]{\linewidth}\raggedright
\begin{quote}
Qwen3-235B
\end{quote}
\end{minipage} & 0.0208 & 0.1316 & 0.0196 & 0.1246 & -0.0951 & -0.60 &
18.56 \\
& \begin{minipage}[t]{\linewidth}\raggedright
\begin{quote}
Deepseek-V3.2
\end{quote}
\end{minipage} & 0.0299 & 0.1969 & 0.0272 & \uline{0.1799} &
\uline{0.3972} & \uline{2.58} & \textbf{9.23} \\
& \begin{minipage}[t]{\linewidth}\raggedright
\begin{quote}
Gemini-3-pro-preview
\end{quote}
\end{minipage} & 0.0263 & 0.1671 & 0.0236 & 0.1512 & 0.1663 & 1.17 &
14.05 \\
& Claude-4.5-sonnet & \uline{0.0311} & \uline{0.2043} & \uline{0.0286} &
0.1754 & \textbf{0.4105} & \textbf{2.84} & 14.72 \\
& \begin{minipage}[t]{\linewidth}\raggedright
\begin{quote}
GPT-5.2
\end{quote}
\end{minipage} & \textbf{0.0347} & \textbf{0.2122} & \textbf{0.0334} &
\textbf{0.2053} & 0.1587 & 1.11 & \uline{13.89} \\
\multirow{2}{*}{} & \begin{minipage}[t]{\linewidth}\raggedright
\begin{quote}
Qwen3-235B
\end{quote}
\end{minipage} & 0.0450 & 0.2538 & 0.0444 & 0.2507 & 0.3511 & 2.06 &
16.36 \\
& \begin{minipage}[t]{\linewidth}\raggedright
\begin{quote}
Deepseek-V3.2
\end{quote}
\end{minipage} & \uline{0.0461} & \uline{0.2624} & \uline{0.0450} &
\uline{0.2574} & \uline{0.6271} & \uline{4.53} & 15.10 \\
\multirow{2}{*}{\begin{minipage}[t]{\linewidth}\raggedright
\begin{minipage}[b]{\linewidth}\raggedright
\textbf{QuantaAlpha}
\end{minipage}
\end{minipage}} & \begin{minipage}[t]{\linewidth}\raggedright
\begin{quote}
Gemini-3-pro-preview
\end{quote}
\end{minipage} & 0.0453 & 0.2551 & 0.0439 & 0.2490 & 0.5834 & 4.21 &
\uline{12.10} \\
& Claude-4.5-sonnet & 0.0445 & 0.2507 & 0.0431 & 0.2446 & 0.5619 & 4.12
& 13.02 \\
& \begin{minipage}[t]{\linewidth}\raggedright
\begin{quote}
GPT-5.2
\end{quote}
\end{minipage} & \begin{minipage}[t]{\linewidth}\raggedright
\begin{minipage}[b]{\linewidth}\raggedright
\textbf{0.0472}
\end{minipage}
\end{minipage} & \begin{minipage}[t]{\linewidth}\raggedright
\begin{minipage}[b]{\linewidth}\raggedright
\textbf{0.2691}
\end{minipage}
\end{minipage} & \begin{minipage}[t]{\linewidth}\raggedright
\begin{minipage}[b]{\linewidth}\raggedright
\textbf{0.0459}
\end{minipage}
\end{minipage} & \begin{minipage}[t]{\linewidth}\raggedright
\begin{minipage}[b]{\linewidth}\raggedright
\textbf{0.2635}
\end{minipage}
\end{minipage} & \begin{minipage}[t]{\linewidth}\raggedright
\begin{minipage}[b]{\linewidth}\raggedright
\textbf{0.6453}
\end{minipage}
\end{minipage} & \begin{minipage}[t]{\linewidth}\raggedright
\begin{minipage}[b]{\linewidth}\raggedright
\textbf{4.68}
\end{minipage}
\end{minipage} & \begin{minipage}[t]{\linewidth}\raggedright
\begin{minipage}[b]{\linewidth}\raggedright
\textbf{11.80}
\end{minipage}
\end{minipage} \\
\bottomrule()
\end{longtable}

\begin{quote}
performance, confirming that each contributes non-trivially to robust
factor generation. The \emph{consistency} gate prevents semantic drift
between the hypothesis, symbolic specification, and implementation. The
\emph{complexity} control improves robustness by discouraging overly
complex expressions that generalize poorly, and its removal leads to the
most pronounced degradation at the strategy level, with the annualized
excess return dropping by 0.95\% and the maximum drawdown increasing by
2.31\%. The \emph{redundancy} control preserves exploration by filtering
near-duplicate structures and mitigating factor crowding. Disabling all
three yields the largest degradation, indicating that these controls are
complementary and jointly needed for reliable generation.

We further analyze cross-seed variance and evolutionary robustness in
Appendix A.5, showing that QuantaAlpha remains stable across different
initial seed combinations and that the evolutionary process does not
drift uncontrollably.

\textbf{5.4} \textbf{More Analysis}

\textbf{Factor Generalizability.} To evaluate factor robustness under
significant market distribution shifts, we conduct a zero-shot transfer
experiment where the factors mined and selected on CSI 300 are directly
deployed on CSI 500 and the S\&P 500 without any re-optimization or
market-specific adaptation. Zero-
\end{quote}

8

\begin{quote}
shot transfer is feasible because the factors use standard OHLCV-based
operators and daily cross-sectional normalization mitigates cross-market
scale differences. As shown in Figure 1, the factors mined on CSI 300
exhibit positive transfer performance on both CSI 500 and the S\&P 500.
On CSI 500, QuantaAlpha remains broadly comparable to other alpha-mining
baselines during the early evaluation period, but shows increasingly
pronounced outperformance as the evaluation horizon extends, ultimately
delivering about 40.28\% cumulative excess return over four years. On the
S\&P 500, QuantaAlpha achieves about 19.1\% cumulative excess return over
the same period, indicating strong robustness under market distribution
shifts.
\end{quote}

\begin{quote}
\textbf{Figure 3:} Ablation study of semantic consistency, complexity, and
redundancy controls.
\end{quote}

\begin{quote}
\includegraphics[width=3.90417in,height=1.80556in]{vertopal_76ea1cc15b3d4460874f5e09e8e3ea79/media/image6.png}

Figure 4: Annual IC and Rank IC comparison on the CSI 300.
\end{quote}

\includegraphics[width=2.26805in,height=1.32778in]{vertopal_76ea1cc15b3d4460874f5e09e8e3ea79/media/image7.png}

\begin{quote}
Figure 5: The evolution of IC over the first five iterations.

The first iteration establishes a baseline factor capturing short-term
momentum, but increased structural complexity coincides with weaker
generalization. Subsequent iterations simplify the expression into a
linear additive form, improving drawdown and stabilizing performance.
The fifth iteration adds participant-differentiated behavioral signals,
incorporating complementary information and further improving
predictability. Throughout the evolution, QuantaAlpha maintains a
cumulative factor pool. Its predictive performance does not saturate
within the first five iterations and only begins to decay after roughly
15 iterations, indicating that the evolution sustains effective
improvement over a substantial horizon before diminishing returns
emerge. Detailed analysis is provided in Appendix E.

\textbf{6 Conclusion}
\end{quote}

We present QuantaAlpha, a self-evolving framework for interpretable
alpha mining that formulates factor discovery as a constrained
multi-agent process. Extensive experiments across Chinese and U.S.
equity markets show that QuantaAlpha consistently produces more stable
and generalizable factors than existing baselines. Beyond empirical
performance, QuantaAlpha improves diversity through broad hypothesis
exploration and redundancy-aware evolution, while maintaining
controllability via symbolic representations and constraint-aware
synthesis. These results suggest that agentic evolution is a promising
paradigm for discovery problems in high-noise, non-stationary domains.

\begin{quote}
\textbf{Impact Statement}

This paper presents work whose goal is to advance the field of Machine
Learning. There are many potential societal consequences of our work,
none which we feel must be specifically highlighted here.

\textbf{References}

Binqi Chen, Hongjun Ding, Ning Shen, Jinsheng Huang, Taian Guo, Luchen
Liu, and Ming Zhang. 2025. Alphasage: Structure-aware alpha mining via
gflownets for robust exploration. \emph{arXiv preprint
arXiv:2509.25055}.

Hongjun Ding, Binqi Chen, Jinsheng Huang, Taian Guo, Zhengyang Mao,
Guoyi Shao, Lutong Zou, Luchen Liu, and Ming Zhang. 2025. Alphaeval: A
comprehensive and efficient evaluation framework for formula alpha
mining. \emph{arXiv preprint arXiv:2508.13174}.

Yitong Duan, chuheng zhang, and Jian Li. 2025. Factormad: A multi-agent
debate framework based on large language models for interpretable stock
alpha factor mining. In \emph{Proceedings of the 6th ACM International
Conference on AI in Finance}, pages 605--613.

Robert F Engle. 1982. Autoregressive conditional heteroscedasticity with
estimates of the variance of united kingdom inflation.
\emph{Econometrica: Journal of the econometric society}, pages
987--1007.

Eugene F Fama. 1965. The behavior of stock-market prices. \emph{The
journal of Business}, 38(1):34--105.

Jinyuan Fang, Yanwen Peng, Xi Zhang, Yingxu Wang, Xinhao Yi, Guibin
Zhang, Yi Xu, Bin Wu, Siwei Liu, Zihao Li, and 1 others. 2025. A
comprehensive survey of self-evolving ai agents: A new paradigm bridging
foundation models and lifelong agentic systems. \emph{arXiv preprint
arXiv:2508.07407}.
\end{quote}

10

\begin{quote}
Xin Guo, Haotian Xia, Zhaowei Liu, Hanyang Cao, Zhi Yang, Zhiqiang Liu,
Sizhe Wang, Jinyi Niu, Chuqi Wang, Yanhui Wang, and 1 others. 2025.
Fineval: A chinese financial domain knowledge evaluation benchmark for
large language models. In \emph{Proceedings of the 2025 Conference of
the Nations of the Americas Chapter of the Association for Computational
Linguistics: Human Language Technologies (Volume 1: Long Papers)}, pages
6258--6292.

Xin Guo, Rongjunchen Zhang, Guilong Lu, Xuntao Guo, Shuai Jia, Zhi Yang,
and Liwen Zhang. 2026. Bizfinbench. v2: A unified dual-mode bilingual
benchmark for expert-level financial capability alignment. \emph{arXiv
preprint arXiv:2601.06401}.

Tu Hu, Ronghao Chen, Shuo Zhang, Jianghao Yin, Mou Xiao Feng, Jingping
Liu, Shaolei Zhang, Wenqi Jiang, Yuqi Fang, Sen Hu, and 1 others. 2026.
Controlled self-evolution for algorithmic code optimization. \emph{arXiv
preprint arXiv:2601.07348}.

Haohang Li, Yupeng Cao, Yangyang Yu, Shashidhar Reddy Javaji, Zhiyang
Deng, Yueru He, Yuechen Jiang, Zining Zhu, Kp Subbalakshmi, Jimin Huang,
and 1 others. 2025a. Investorbench: A benchmark for financial
decision-making tasks with llm-based agent. In \emph{Proceedings of the
63rd Annual Meeting of the Association for Computational Linguistics
(Volume 1: Long Papers)}, pages 2509--2525.

Xiang Li, Zhenyu Li, Chen Shi, Yong Xu, Qing Du, Mingkui Tan, and Jun
Huang. 2024a. Alphafin: Benchmarking financial analysis with
retrieval-augmented stock-chain framework. In \emph{Proceedings of the
2024 joint international conference on computational linguistics,
language resources and evaluation (LREC-COLING 2024)}, pages 773--783.

Xiangyu Li, Yawen Zeng, Xiaofen Xing, Jin Xu, and Xiangmin Xu. 2025b.
Hedgeagents: A balanced-aware multi-agent financial trading system. In
\emph{Companion Proceedings of the ACM on Web Conference 2025}, pages
296--305.

Xiangyu Li, Yawen Zeng, Xiaofen Xing, Jin Xu, and Xiangmin Xu. 2025c.
Quantagents: Towards multi-agent financial system via simulated trading.
\emph{arXiv preprint arXiv:2510.04643}.

Yuante Li, Xu Yang, Xiao Yang, Minrui Xu, Xisen Wang, Weiqing Liu, and
Jiang Bian. 2025d. R\&d-agent-quant: a multi-agent framework for
data-centric factors and model joint optimization. \emph{arXiv preprint
arXiv:2505.15155}.

Zhiwei Li, Ran Song, Caihong Sun, Wei Xu, Zhengtao Yu, and Ji-Rong Wen.
2024b. Can large language models mine interpretable financial factors
more effectively? a neural-symbolic factor mining agent model. In
\emph{Findings of the Association for Computational Linguistics ACL
2024}, pages 3891--3902.

Jiaye Lin, Yifu Guo, Yuzhen Han, Sen Hu, Ziyi Ni, Licheng Wang,
Mingguang Chen, Hongzhang Liu, Ronghao Chen, Yangfan He, and 1 others.
2025. Se-agent: Self-evolution trajectory optimization in multi-step
reasoning with llm-based agents. \emph{arXiv preprint arXiv:2508.02085}.

Fengyuan Liu, Huang Yi, Sichun Luo, Yuqi Wang, Yazheng Yang, Xinye Li,
Zefa Hu, Junlan Feng, and Qi Liu. 2025a. Cognitive alpha mining via
llm-driven code-based evolution. \emph{arXiv preprint arXiv:2511.18850}.

Xiao-Yang Liu, Guoxuan Wang, Hongyang Yang, and Daochen Zha. 2023.
Fingpt: Democratizing internet-scale data for financial large language
models. \emph{arXiv preprint arXiv:2307.10485}.

Zhaowei Liu, Xin Guo, Zhi Yang, Fangqi Lou, Lingfeng Zeng, Jinyi Niu,
Mengping Li, Qi Qi, Zhiqiang Liu, Yiyang Han, and 1 others. 2025b.
Fin-r1: A large language model for financial reasoning through
reinforcement learning. \emph{arXiv preprint arXiv:2503.16252}.

Alejandro Lopez-Lira and Yuehua Tang. 2023. Can chatgpt forecast stock
price movements? return predictability and large language models.
\emph{arXiv preprint arXiv:2304.07619}.

Alexander Novikov, Ngân V˜u, Marvin Eisenberger, Emilien Dupont, Po-Sen
Huang, Adam Zsolt Wagner, Sergey Shirobokov, Borislav Kozlovskii,
Francisco JR Ruiz, Abbas Mehrabian, and 1 others. 2025. Alphaevolve: A
coding agent for scientific and algorithmic discovery. \emph{arXiv
preprint arXiv:2506.13131}.

Charidimos Papadakis, Angeliki Dimitriou, Giorgos Filandrianos, Maria
Lymperaiou, Konstantinos Thomas, and Giorgos Stamou. 2025. Atlas:
Adaptive trading with llm agents through dynamic prompt optimization and
multi-agent coordination. \emph{arXiv preprint arXiv:2510.15949}.

M Hashem Pesaran. 2021. General diagnostic tests for cross-sectional
dependence in panels. \emph{Empirical economics}, 60(1):13--50.
\end{quote}

11

\begin{quote}
Hao Shi, Weili Song, Xinting Zhang, Jiahe Shi, Cuicui Luo, Xiang Ao,
Hamid Arian, and Luis Angel Seco. 2025a. Alphaforge: A framework to mine
and dynamically combine formulaic alpha factors. In \emph{Proceedings of
the AAAI conference on artificial intelligence}, volume 39, pages
12524--12532.

Yu Shi, Yitong Duan, and Jian Li. 2025b. Navigating the alpha jungle: An
llm-powered mcts framework for formulaic factor mining. \emph{arXiv
preprint arXiv:2505.11122}.

Zichen Tang, Jiacheng Liu, Zhongjun Yang, Rongjin Li, Zihua Rong,
Haoyang He, Zhuodi Hao, Xinyang Hu, Kun Ji, Ziyan Ma, and 1 others.
2025a. Finmmr: make financial numerical reasoning more multimodal,
comprehensive, and challenging. In \emph{Proceedings of the IEEE/CVF
International Conference on Computer Vision}, pages 3245--3257.

Ziyi Tang, Zechuan Chen, Jiarui Yang, Jiayao Mai, Yongsen Zheng, Keze
Wang, Jinrui Chen, and Liang Lin. 2025b. Alphaagent: Llm-driven alpha
mining with regularized exploration to counteract alpha decay. In
\emph{Proceedings of the 31st ACM SIGKDD Conference on Knowledge
Discovery and Data Mining V. 2}, pages 2813--2822.

Saizhuo Wang, Hang Yuan, Leon Zhou, Lionel Ni, Heung Yeung Shum, and
Jian Guo. 2025. Alpha-gpt: Human-ai interactive alpha mining for
quantitative investment. In \emph{Proceedings of the 2025 Conference on
Empirical Methods in Natural Language Processing: System
Demonstrations}, pages 196--206.

Yijia Xiao, Edward Sun, Di Luo, and Wei Wang. 2024. Tradingagents:
Multi-agents llm financial trading framework. \emph{arXiv preprint
arXiv:2412.20138}.

Zhi Yang, Runguo Li, Qiqi Qiang, Jiashun Wang, Fangqi Lou, Mengping Li,
Dongpo Cheng, Rui Xu, Heng Lian, Shuo Zhang, and 1 others. 2026.
Finvault: Benchmarking financial agent safety in execution-grounded
environments. \emph{arXiv preprint arXiv:2601.07853}.

Yangyang Yu, Haohang Li, Zhi Chen, Yuechen Jiang, Yang Li, Jordan W
Suchow, Denghui Zhang, and Khaldoun Khashanah. 2025. Finmem: A
performance-enhanced llm trading agent with layered memory and character
design. \emph{IEEE Transactions on Big Data}.

Yangyang Yu, Zhiyuan Yao, Haohang Li, Zhiyang Deng, Yuechen Jiang,
Yupeng Cao, Zhi Chen, Jordan W Suchow, Zhenyu Cui, Rong Liu, and 1
others. 2024. Fincon: A synthesized llm multi-agent system with
conceptual verbal reinforcement for enhanced financial decision making.
\emph{Advances in Neural Information Processing Systems},
37:137010--137045.

Yunpeng Zhai, Shuchang Tao, Cheng Chen, Anni Zou, Ziqian Chen, Qingxu
Fu, Shinji Mai, Li Yu, Jiaji Deng, Zouying Cao, and 1 others. 2025.
Agentevolver: Towards efficient self-evolving agent system. \emph{arXiv
preprint arXiv:2511.10395}.

Shuo Zhang, Chaofa Yuan, Ryan Guo, Xiaomin Yu, Rui Xu, Zhangquan Chen,
Zinuo Li, Zhi Yang, Shuhao Guan, Zhenheng Tang, and 1 others. 2026.
Evofsm: Controllable self-evolution for deep research with finite state
machines. \emph{arXiv preprint arXiv:2601.09465}.

Wentao Zhang, Lingxuan Zhao, Haochong Xia, Shuo Sun, Jiaze Sun, Molei
Qin, Xinyi Li, Yuqing Zhao, Yilei Zhao, Xinyu Cai, and 1 others. 2024. A
multimodal foundation agent for financial trading: Tool-augmented,
diversified, and generalist. In \emph{Proceedings of the 30th acm sigkdd
conference on knowledge discovery and data mining}, pages 4314--4325.
\end{quote}

12

\begin{quote}
\textbf{SUMMARY OF THE APPENDIX}\\
This appendix contains additional details for the
\emph{\textbf{``QuantaAlpha: An Evolutionary Framework for LLM-Driven
Alpha Mining''}}. The appendix is organized as follows:

\textbf{Appendix Table of Contents}
\end{quote}

\begin{longtable}[]{@{}
  >{\raggedright\arraybackslash}p{(\columnwidth - 16\tabcolsep) * \real{0.1111}}
  >{\raggedright\arraybackslash}p{(\columnwidth - 16\tabcolsep) * \real{0.1111}}
  >{\raggedright\arraybackslash}p{(\columnwidth - 16\tabcolsep) * \real{0.1111}}
  >{\raggedright\arraybackslash}p{(\columnwidth - 16\tabcolsep) * \real{0.1111}}
  >{\raggedright\arraybackslash}p{(\columnwidth - 16\tabcolsep) * \real{0.1111}}
  >{\raggedright\arraybackslash}p{(\columnwidth - 16\tabcolsep) * \real{0.1111}}
  >{\raggedright\arraybackslash}p{(\columnwidth - 16\tabcolsep) * \real{0.1111}}
  >{\raggedright\arraybackslash}p{(\columnwidth - 16\tabcolsep) * \real{0.1111}}
  >{\raggedright\arraybackslash}p{(\columnwidth - 16\tabcolsep) * \real{0.1111}}@{}}
\toprule()
\multicolumn{8}{@{}>{\raggedright\arraybackslash}p{(\columnwidth - 16\tabcolsep) * \real{0.8889} + 14\tabcolsep}}{%
\begin{minipage}[b]{\linewidth}\raggedright
\begin{quote}
\textbf{A Experiment Settings}
\end{quote}
\end{minipage}} &
\multirow{22}{*}{\begin{minipage}[b]{\linewidth}\raggedright
\begin{quote}
\textbf{13} 13 14 14 14 15 16
\end{quote}

\textbf{16}

\begin{quote}
\textbf{16} 17 18 18 19 19

\textbf{20} 20 20 21 22 22 22
\end{quote}

\textbf{22}
\end{minipage}} \\
\multicolumn{2}{@{}>{\raggedright\arraybackslash}p{(\columnwidth - 16\tabcolsep) * \real{0.2222} + 2\tabcolsep}}{%
\begin{minipage}[b]{\linewidth}\raggedright
A.1
\end{minipage}} &
\multicolumn{2}{>{\raggedright\arraybackslash}p{(\columnwidth - 16\tabcolsep) * \real{0.2222} + 2\tabcolsep}}{%
\begin{minipage}[b]{\linewidth}\raggedright
Evaluation Metrics
\end{minipage}} &
\multicolumn{4}{>{\raggedright\arraybackslash}p{(\columnwidth - 16\tabcolsep) * \real{0.4444} + 6\tabcolsep}}{%
\begin{minipage}[b]{\linewidth}\raggedright
\begin{quote}
. . . . . . . . . . . . . . . . . . . . . . . . . . . . . . . .
\end{quote}
\end{minipage}} \\
\multicolumn{2}{@{}>{\raggedright\arraybackslash}p{(\columnwidth - 16\tabcolsep) * \real{0.2222} + 2\tabcolsep}}{%
\begin{minipage}[b]{\linewidth}\raggedright
A.2
\end{minipage}} &
\multicolumn{6}{>{\raggedright\arraybackslash}p{(\columnwidth - 16\tabcolsep) * \real{0.6667} + 10\tabcolsep}}{%
\begin{minipage}[b]{\linewidth}\raggedright
Backtesting Setup . . . . . . . . . . . . . . . . . . . . . . . . . . .
. . . . . .
\end{minipage}} \\
\multicolumn{2}{@{}>{\raggedright\arraybackslash}p{(\columnwidth - 16\tabcolsep) * \real{0.2222} + 2\tabcolsep}}{%
\begin{minipage}[b]{\linewidth}\raggedright
A.3
\end{minipage}} & \begin{minipage}[b]{\linewidth}\raggedright
\begin{quote}
Baselines
\end{quote}
\end{minipage} &
\multicolumn{5}{>{\raggedright\arraybackslash}p{(\columnwidth - 16\tabcolsep) * \real{0.5556} + 8\tabcolsep}}{%
\begin{minipage}[b]{\linewidth}\raggedright
. . . . . . . . . . . . . . . . . . . . . . . . . . . . . . . . . . . .
.
\end{minipage}} \\
\multicolumn{2}{@{}>{\raggedright\arraybackslash}p{(\columnwidth - 16\tabcolsep) * \real{0.2222} + 2\tabcolsep}}{%
\begin{minipage}[b]{\linewidth}\raggedright
A.4
\end{minipage}} &
\multicolumn{4}{>{\raggedright\arraybackslash}p{(\columnwidth - 16\tabcolsep) * \real{0.4444} + 6\tabcolsep}}{%
\begin{minipage}[b]{\linewidth}\raggedright
\begin{quote}
Computational Cost and Token Consumption
\end{quote}
\end{minipage}} &
\multicolumn{2}{>{\raggedright\arraybackslash}p{(\columnwidth - 16\tabcolsep) * \real{0.2222} + 2\tabcolsep}}{%
\begin{minipage}[b]{\linewidth}\raggedright
. . . . . . . . . . . . . . . . . .
\end{minipage}} \\
\multicolumn{2}{@{}>{\raggedright\arraybackslash}p{(\columnwidth - 16\tabcolsep) * \real{0.2222} + 2\tabcolsep}}{%
\begin{minipage}[b]{\linewidth}\raggedright
A.5
\end{minipage}} &
\multicolumn{5}{>{\raggedright\arraybackslash}p{(\columnwidth - 16\tabcolsep) * \real{0.5556} + 8\tabcolsep}}{%
\begin{minipage}[b]{\linewidth}\raggedright
Cross-Seed Variance and Evolutionary Robustness
\end{minipage}} & \begin{minipage}[b]{\linewidth}\raggedright
\begin{quote}
. . . . . . . . . . . . . . .
\end{quote}
\end{minipage} \\
\multicolumn{5}{@{}>{\raggedright\arraybackslash}p{(\columnwidth - 16\tabcolsep) * \real{0.5556} + 8\tabcolsep}}{%
\begin{minipage}[b]{\linewidth}\raggedright
\begin{quote}
A.6 Transaction Cost Sensitivity Analysis
\end{quote}
\end{minipage}} &
\multicolumn{3}{>{\raggedright\arraybackslash}p{(\columnwidth - 16\tabcolsep) * \real{0.3333} + 4\tabcolsep}}{%
\begin{minipage}[b]{\linewidth}\raggedright
. . . . . . . . . . . . . . . . . . . . . .
\end{minipage}} \\
\multicolumn{8}{@{}>{\raggedright\arraybackslash}p{(\columnwidth - 16\tabcolsep) * \real{0.8889} + 14\tabcolsep}}{%
\begin{minipage}[b]{\linewidth}\raggedright
\begin{quote}
\textbf{B Algorithm Configuration}\\
\textbf{C Case Study: Factor Evolution Trajectory}
\end{quote}\strut
\end{minipage}} \\
\multicolumn{2}{@{}>{\raggedright\arraybackslash}p{(\columnwidth - 16\tabcolsep) * \real{0.2222} + 2\tabcolsep}}{%
\begin{minipage}[b]{\linewidth}\raggedright
C.1
\end{minipage}} &
\multicolumn{6}{>{\raggedright\arraybackslash}p{(\columnwidth - 16\tabcolsep) * \real{0.6667} + 10\tabcolsep}}{%
\begin{minipage}[b]{\linewidth}\raggedright
Factor Identity . . . . . . . . . . . . . . . . . . . . . . . . . . . .
. . . . . . .
\end{minipage}} \\
\multicolumn{2}{@{}>{\raggedright\arraybackslash}p{(\columnwidth - 16\tabcolsep) * \real{0.2222} + 2\tabcolsep}}{%
\begin{minipage}[b]{\linewidth}\raggedright
C.2
\end{minipage}} &
\multicolumn{6}{>{\raggedright\arraybackslash}p{(\columnwidth - 16\tabcolsep) * \real{0.6667} + 10\tabcolsep}}{%
\begin{minipage}[b]{\linewidth}\raggedright
Evolution Lineage . . . . . . . . . . . . . . . . . . . . . . . . . . .
. . . . . .
\end{minipage}} \\
\multicolumn{2}{@{}>{\raggedright\arraybackslash}p{(\columnwidth - 16\tabcolsep) * \real{0.2222} + 2\tabcolsep}}{%
\begin{minipage}[b]{\linewidth}\raggedright
C.3
\end{minipage}} &
\multicolumn{6}{>{\raggedright\arraybackslash}p{(\columnwidth - 16\tabcolsep) * \real{0.6667} + 10\tabcolsep}}{%
\begin{minipage}[b]{\linewidth}\raggedright
Synthesized Hypothesis . . . . . . . . . . . . . . . . . . . . . . . . .
. . . . .
\end{minipage}} \\
\multicolumn{2}{@{}>{\raggedright\arraybackslash}p{(\columnwidth - 16\tabcolsep) * \real{0.2222} + 2\tabcolsep}}{%
\begin{minipage}[b]{\linewidth}\raggedright
C.4
\end{minipage}} &
\multicolumn{6}{>{\raggedright\arraybackslash}p{(\columnwidth - 16\tabcolsep) * \real{0.6667} + 10\tabcolsep}}{%
\begin{minipage}[b]{\linewidth}\raggedright
Backtest Performance . . . . . . . . . . . . . . . . . . . . . . . . . .
. . . . .
\end{minipage}} \\
\multicolumn{2}{@{}>{\raggedright\arraybackslash}p{(\columnwidth - 16\tabcolsep) * \real{0.2222} + 2\tabcolsep}}{%
\begin{minipage}[b]{\linewidth}\raggedright
C.5
\end{minipage}} &
\multicolumn{2}{>{\raggedright\arraybackslash}p{(\columnwidth - 16\tabcolsep) * \real{0.2222} + 2\tabcolsep}}{%
\begin{minipage}[b]{\linewidth}\raggedright
\begin{quote}
Trajectory Summary
\end{quote}
\end{minipage}} &
\multicolumn{4}{>{\raggedright\arraybackslash}p{(\columnwidth - 16\tabcolsep) * \real{0.4444} + 6\tabcolsep}}{%
\begin{minipage}[b]{\linewidth}\raggedright
. . . . . . . . . . . . . . . . . . . . . . . . . . . . . . .
\end{minipage}} \\
\multicolumn{8}{@{}>{\raggedright\arraybackslash}p{(\columnwidth - 16\tabcolsep) * \real{0.8889} + 14\tabcolsep}}{%
\begin{minipage}[b]{\linewidth}\raggedright
\begin{quote}
\textbf{D Market Style, Factor Semantics, and Alpha Decay}
\end{quote}
\end{minipage}} \\
\multicolumn{2}{@{}>{\raggedright\arraybackslash}p{(\columnwidth - 16\tabcolsep) * \real{0.2222} + 2\tabcolsep}}{%
\begin{minipage}[b]{\linewidth}\raggedright
D.1
\end{minipage}} &
\multicolumn{6}{>{\raggedright\arraybackslash}p{(\columnwidth - 16\tabcolsep) * \real{0.6667} + 10\tabcolsep}}{%
\begin{minipage}[b]{\linewidth}\raggedright
Market style transition and its implications . . . . . . . . . . . . . .
. . . . . .
\end{minipage}} \\
\multicolumn{2}{@{}>{\raggedright\arraybackslash}p{(\columnwidth - 16\tabcolsep) * \real{0.2222} + 2\tabcolsep}}{%
\begin{minipage}[b]{\linewidth}\raggedright
D.2
\end{minipage}} &
\multicolumn{5}{>{\raggedright\arraybackslash}p{(\columnwidth - 16\tabcolsep) * \real{0.5556} + 8\tabcolsep}}{%
\begin{minipage}[b]{\linewidth}\raggedright
\begin{quote}
Factor semantics aligned with market microstructure
\end{quote}
\end{minipage}} & \begin{minipage}[b]{\linewidth}\raggedright
\begin{quote}
. . . . . . . . . . . . . .
\end{quote}
\end{minipage} \\
\multicolumn{2}{@{}>{\raggedright\arraybackslash}p{(\columnwidth - 16\tabcolsep) * \real{0.2222} + 2\tabcolsep}}{%
\begin{minipage}[b]{\linewidth}\raggedright
D.3
\end{minipage}} &
\multicolumn{6}{>{\raggedright\arraybackslash}p{(\columnwidth - 16\tabcolsep) * \real{0.6667} + 10\tabcolsep}}{%
\begin{minipage}[b]{\linewidth}\raggedright
Factor statistics as supporting evidence . . . . . . . . . . . . . . . .
. . . . . .
\end{minipage}} \\
\multicolumn{2}{@{}>{\raggedright\arraybackslash}p{(\columnwidth - 16\tabcolsep) * \real{0.2222} + 2\tabcolsep}}{%
\begin{minipage}[b]{\linewidth}\raggedright
D.4
\end{minipage}} &
\multicolumn{6}{>{\raggedright\arraybackslash}p{(\columnwidth - 16\tabcolsep) * \real{0.6667} + 10\tabcolsep}}{%
\begin{minipage}[b]{\linewidth}\raggedright
Role of factor diversity and mutation . . . . . . . . . . . . . . . . .
. . . . . .
\end{minipage}} \\
\multicolumn{2}{@{}>{\raggedright\arraybackslash}p{(\columnwidth - 16\tabcolsep) * \real{0.2222} + 2\tabcolsep}}{%
\begin{minipage}[b]{\linewidth}\raggedright
D.5
\end{minipage}} & \begin{minipage}[b]{\linewidth}\raggedright
\begin{quote}
Summary
\end{quote}
\end{minipage} &
\multicolumn{5}{>{\raggedright\arraybackslash}p{(\columnwidth - 16\tabcolsep) * \real{0.5556} + 8\tabcolsep}}{%
\begin{minipage}[b]{\linewidth}\raggedright
. . . . . . . . . . . . . . . . . . . . . . . . . . . . . . . . . . . .
.
\end{minipage}} \\
\multicolumn{2}{@{}>{\raggedright\arraybackslash}p{(\columnwidth - 16\tabcolsep) * \real{0.2222} + 2\tabcolsep}}{%
\begin{minipage}[b]{\linewidth}\raggedright
D.6
\end{minipage}} &
\multicolumn{6}{>{\raggedright\arraybackslash}p{(\columnwidth - 16\tabcolsep) * \real{0.6667} + 10\tabcolsep}}{%
\begin{minipage}[b]{\linewidth}\raggedright
Stress-Period Analysis (High-Volatility Episode in April 2025) . . . . .
. . . .
\end{minipage}} \\
\begin{minipage}[b]{\linewidth}\raggedright
\textbf{E}
\end{minipage} &
\multicolumn{7}{>{\raggedright\arraybackslash}p{(\columnwidth - 16\tabcolsep) * \real{0.7778} + 12\tabcolsep}}{%
\begin{minipage}[b]{\linewidth}\raggedright
\begin{quote}
\textbf{Iteration Study: Convergence of Factor-Pool Performance
(CSI300)}
\end{quote}
\end{minipage}} \\
\begin{minipage}[b]{\linewidth}\raggedright
\textbf{A}
\end{minipage} &
\multicolumn{7}{>{\raggedright\arraybackslash}p{(\columnwidth - 16\tabcolsep) * \real{0.7778} + 12\tabcolsep}}{%
\begin{minipage}[b]{\linewidth}\raggedright
\begin{quote}
\textbf{Experiment Settings}
\end{quote}
\end{minipage}} \\
\midrule()
\endhead
\bottomrule()
\end{longtable}

\begin{quote}
This section provides details of our experimental setup, including
computational infrastructure, evalu-ation metrics, backtesting setup,
and baselines. Data and code are available at.

\textbf{A.1 Evaluation Metrics}
\end{quote}

We evaluate predictive performance using two categories of metrics:
factor predictive power and strategy-level performance.

\begin{quote}
Without loss of generality, the bar notation¯\emph{·} denotes the mean,
and \emph{$\sigma$}(\emph{·}) denotes the standard deviation.

• \textbf{Information Coefficient (IC)}: Pearson correlation between
factor values \textbf{f}\emph{t} and future returns
\textbf{r}\emph{t}+1:
\end{quote}

\begin{longtable}[]{@{}
  >{\raggedright\arraybackslash}p{(\columnwidth - 4\tabcolsep) * \real{0.3333}}
  >{\raggedright\arraybackslash}p{(\columnwidth - 4\tabcolsep) * \real{0.3333}}
  >{\raggedright\arraybackslash}p{(\columnwidth - 4\tabcolsep) * \real{0.3333}}@{}}
\toprule()
\begin{minipage}[b]{\linewidth}\raggedright
IC\emph{t} =
\end{minipage} & \begin{minipage}[b]{\linewidth}\raggedright
(\textbf{f}\emph{t $-$}¯\emph{ft}\textbf{1})\emph{${}^{\top}$}(\textbf{r}\emph{t}+1
\emph{$-$}¯\emph{rt}+1\textbf{1})

\emph{$\|$}\textbf{f}\emph{t $-$}¯\emph{ft}\textbf{1}\emph{$\|$}2 \emph{·
$\|$}\textbf{r}\emph{t}+1 \emph{$-$}¯\emph{rt}+1\textbf{1}\emph{$\|$}2
\end{minipage} & \begin{minipage}[b]{\linewidth}\raggedright
\begin{quote}
\emph{,}
\end{quote}
\end{minipage} \\
\midrule()
\endhead
\bottomrule()
\end{longtable}

\begin{quote}
where \textbf{1} denotes a column vector of ones.

• \textbf{ICIR}: Information ratio of IC, measuring consistency: ICIR =
IC\emph{/$\sigma$}(IC).
\end{quote}

13

\begin{longtable}[]{@{}
  >{\raggedright\arraybackslash}p{(\columnwidth - 4\tabcolsep) * \real{0.3333}}
  >{\raggedright\arraybackslash}p{(\columnwidth - 4\tabcolsep) * \real{0.3333}}
  >{\raggedright\arraybackslash}p{(\columnwidth - 4\tabcolsep) * \real{0.3333}}@{}}
\toprule()
\multicolumn{3}{@{}>{\raggedright\arraybackslash}p{(\columnwidth - 4\tabcolsep) * \real{1.0000} + 4\tabcolsep}@{}}{%
\begin{minipage}[b]{\linewidth}\raggedright
\begin{quote}
• \textbf{Rank IC}: Spearman correlation using rank
vectors˜\textbf{f}\emph{t} = rank(\textbf{f}\emph{t}) and
˜\textbf{r}\emph{t}+1 = rank(\textbf{r}\emph{t}+1):
\end{quote}
\end{minipage}} \\
\midrule()
\endhead
RankIC\emph{t} = & \begin{minipage}[t]{\linewidth}\raggedright
\begin{quote}
(˜\textbf{f}\emph{t $-$}¯ \textbf{r}\emph{t}+1 \emph{$-$}¯

\emph{$\|$}˜\textbf{f}\emph{t $-$}¯\emph{ft}\textbf{1}\emph{$\|$}2 \emph{·
$\|$}˜\textbf{r}\emph{t}+1 \emph{$-$}¯
\end{quote}
\end{minipage} & \begin{minipage}[t]{\linewidth}\raggedright
\begin{quote}
\emph{,}
\end{quote}
\end{minipage} \\
\bottomrule()
\end{longtable}

\begin{quote}
where rank(\emph{·}) is the rank function applied element-wise to its
input vector in ascending order.• \textbf{Rank ICIR}: Information ratio
of Rank IC: RankICIR = RankIC\emph{/$\sigma$}(RankIC).
\end{quote}

\begin{longtable}[]{@{}
  >{\raggedright\arraybackslash}p{(\columnwidth - 0\tabcolsep) * \real{1.0000}}@{}}
\toprule()
\begin{minipage}[b]{\linewidth}\raggedright
\begin{quote}
All strategy metrics are computed on \emph{excess returns after
transaction costs}, where \emph{r}excess\emph{,t} =
\emph{r}portfolio\emph{,t$-$r}benchmark\emph{,t $-$c}transaction\emph{,t}.
Here, \emph{r}portfolio\emph{,t} represents the return of the strategy
portfolio, \emph{r}benchmark\emph{,t} denotes the return of the market
benchmark, and \emph{c}transaction\emph{,t} accounts for the transaction
costs incurred at time \emph{t}.

• \textbf{Information Ratio (}IR\textbf{)}: IR =
(\emph{r}excess\emph{/$\sigma$}(\emph{r}excess)) \emph{×}$\sqrt{252}$.

• \textbf{Annualized Return (}ARR\textbf{)}: Annualized excess return
over benchmark.
\end{quote}
\end{minipage} \\
\midrule()
\endhead
\bottomrule()
\end{longtable}

\begin{quote}
• \textbf{Maximum Drawdown (}MDD\textbf{)}: Largest peak-to-trough
decline in cumulative excess returns.

\textbf{A.2 Backtesting Setup}

Backtesting is conducted using the Qlib framework across the CSI 300,
CSI 500, and S\&P 500 indices, with the data split detailed in Table 3.
Factor construction utilizes six basic features---open, high, low,
close, volume, and vwap---to predict the next-day return, defined as
\emph{yt} = \emph{P}close denotes the closing price at time \emph{t}. To
ensure robustness against outliers, the preprocessing pipeline
\emph{t}+2\emph{/P} close \emph{t}+1\emph{$-$}1, where \emph{P} close

includes forward-filling missing values, replacing infinite values,
dropping samples with missing labels, and applying cross-sectional rank
normalization (CSRankNorm) to both features and labels.
\end{quote}

Table 3: Data Split Periods for Train, Validation, and Test Sets across
All Markets

\begin{longtable}[]{@{}
  >{\raggedright\arraybackslash}p{(\columnwidth - 8\tabcolsep) * \real{0.2000}}
  >{\raggedright\arraybackslash}p{(\columnwidth - 8\tabcolsep) * \real{0.2000}}
  >{\raggedright\arraybackslash}p{(\columnwidth - 8\tabcolsep) * \real{0.2000}}
  >{\raggedright\arraybackslash}p{(\columnwidth - 8\tabcolsep) * \real{0.2000}}
  >{\raggedright\arraybackslash}p{(\columnwidth - 8\tabcolsep) * \real{0.2000}}@{}}
\toprule()
\begin{minipage}[b]{\linewidth}\raggedright
\end{minipage} & \begin{minipage}[b]{\linewidth}\raggedright
\begin{quote}
\textbf{Market}
\end{quote}
\end{minipage} & \begin{minipage}[b]{\linewidth}\raggedright
\textbf{Train}
\end{minipage} & \begin{minipage}[b]{\linewidth}\raggedright
\textbf{Valid}
\end{minipage} & \begin{minipage}[b]{\linewidth}\raggedright
\textbf{Test}
\end{minipage} \\
\midrule()
\endhead
\multirow{3}{*}{} & CSI 300 & 2016-01-01--2020-12-31 &
2021-01-01--2021-12-31 & 2022-01-01--2025-12-26 \\
& CSI 500 & 2016-01-01--2020-12-31 & 2021-01-01--2021-12-31 &
2022-01-01--2025-12-26 \\
& S\&P 500 & 2016-01-01--2020-12-31 & 2021-01-01--2021-12-31 &
2022-01-01--2025-12-26 \\
\begin{minipage}[t]{\linewidth}\raggedright
\begin{quote}
\textbf{A.3}
\end{quote}
\end{minipage} & \textbf{Baselines} & & & \\
\bottomrule()
\end{longtable}

We benchmark against four categories: (1) \textbf{ML models}: Linear
Regression (Linear), Multi-Layer Perceptron (MLP), and gradient boosting
decision trees including LightGBM, XGBoost, and CatBoost, along with
DoubleEnsemble, an ensemble method for financial time series; (2)
\textbf{Deep learning}: Recurrent networks such as Gated Recurrent Unit
(GRU) and Long Short-Term Memory (LSTM), the attention-based
Transformer, and Temporal Routing Adaptor (TRA); (3) \textbf{Classical
factors}: Alpha158 and Alpha360, which are widely used sets of technical
factors derived from price and volume; (4) \textbf{LLM agents}: RD-Agent
and AlphaAgent, which utilize large language models for automated factor
mining. For the LLM-agent baselines (RD-Agent and AlphaAgent), we
evaluate multiple backbone LLMs, including Qwen3-235B, DeepSeek-V3.2,
Gemini-3-Pro-Preview, Claude-4.5-Sonnet, and GPT-5.2. Unless otherwise
specified, all other LLM-based experiments in this paper adopt
DeepSeek-V3.2 for consistency and fair comparison.

\begin{quote}
\textbf{A.4 Computational Cost and Token Consumption}

Under the standard main-experiment setting, QuantaAlpha uses 10 parallel
planning directions and 5 main evolutionary iterations, consisting of
one initial round followed by mutation and crossover rounds. A complete
run takes approximately 20 hours using hosted LLM APIs and CPU-based
factor evaluation/backtesting, without requiring local GPU resources.

In terms of LLM usage, QuantaAlpha consumes approximately 1.8M tokens
per main run on average, compared with 1.5M tokens for AlphaAgent and
2.3M tokens for RD-Agent under comparable generated factor scales. The
additional cost over AlphaAgent mainly comes from maintaining
trajectory-pool
\end{quote}

14

\begin{quote}
information, intermediate-state records, consistency verification, and
iterative feedback for mutation and crossover. Compared with RD-Agent,
QuantaAlpha achieves stronger performance with a lower token budget,
suggesting a favorable cost-performance trade-off. Once mined, the
resulting symbolic factor pool can be reused for subsequent backtesting,
portfolio construction, and cross-market transfer without rerunning the
full evolutionary search.

\textbf{A.5 Cross-Seed Variance and Evolutionary Robustness}
\end{quote}

To assess whether QuantaAlpha is sensitive to initialization randomness,
we conduct a cross-seed variance analysis using three representative
initial seed combinations. The initial seed factor pool is derived from
the public Alpha158(20) subset and grouped by low correlation to form
different starting combinations. Table 4 reports the main predictive
metrics under different seed combinations, and Table 5 summarizes the
corresponding variance statistics.

Table 4: Core metrics under different initial seed combinations.

\begin{longtable}[]{@{}
  >{\raggedright\arraybackslash}p{(\columnwidth - 8\tabcolsep) * \real{0.2000}}
  >{\raggedright\arraybackslash}p{(\columnwidth - 8\tabcolsep) * \real{0.2000}}
  >{\raggedright\arraybackslash}p{(\columnwidth - 8\tabcolsep) * \real{0.2000}}
  >{\raggedright\arraybackslash}p{(\columnwidth - 8\tabcolsep) * \real{0.2000}}
  >{\raggedright\arraybackslash}p{(\columnwidth - 8\tabcolsep) * \real{0.2000}}@{}}
\toprule()
\begin{minipage}[b]{\linewidth}\raggedright
Seed Combination
\end{minipage} & \begin{minipage}[b]{\linewidth}\raggedright
IC
\end{minipage} & \begin{minipage}[b]{\linewidth}\raggedright
ICIR
\end{minipage} & \begin{minipage}[b]{\linewidth}\raggedright
Rank IC
\end{minipage} & \begin{minipage}[b]{\linewidth}\raggedright
Rank ICIR
\end{minipage} \\
\midrule()
\endhead
\begin{minipage}[t]{\linewidth}\raggedright
\begin{quote}
Combination 1
\end{quote}
\end{minipage} & 0.0466 & 0.2708 & 0.0454 & 0.2655 \\
\begin{minipage}[t]{\linewidth}\raggedright
\begin{quote}
Combination 2
\end{quote}
\end{minipage} & 0.0426 & 0.2325 & 0.0409 & 0.2236 \\
\begin{minipage}[t]{\linewidth}\raggedright
\begin{quote}
Combination 3
\end{quote}
\end{minipage} & 0.0436 & 0.2551 & 0.0418 & 0.2468 \\
\bottomrule()
\end{longtable}

Table 5: Cross-seed variance summary.

\begin{longtable}[]{@{}
  >{\raggedright\arraybackslash}p{(\columnwidth - 8\tabcolsep) * \real{0.2000}}
  >{\raggedright\arraybackslash}p{(\columnwidth - 8\tabcolsep) * \real{0.2000}}
  >{\raggedright\arraybackslash}p{(\columnwidth - 8\tabcolsep) * \real{0.2000}}
  >{\raggedright\arraybackslash}p{(\columnwidth - 8\tabcolsep) * \real{0.2000}}
  >{\raggedright\arraybackslash}p{(\columnwidth - 8\tabcolsep) * \real{0.2000}}@{}}
\toprule()
\begin{minipage}[b]{\linewidth}\raggedright
\begin{quote}
Metric
\end{quote}
\end{minipage} & \begin{minipage}[b]{\linewidth}\raggedright
Mean
\end{minipage} & \begin{minipage}[b]{\linewidth}\raggedright
Std Dev
\end{minipage} & \begin{minipage}[b]{\linewidth}\raggedright
Coefficient of Variation
\end{minipage} & \begin{minipage}[b]{\linewidth}\raggedright
Range
\end{minipage} \\
\midrule()
\endhead
\begin{minipage}[t]{\linewidth}\raggedright
\begin{quote}
IC
\end{quote}
\end{minipage} & 0.0443 & 0.0021 & 4.64\% & 0.0040 \\
\begin{minipage}[t]{\linewidth}\raggedright
\begin{quote}
ICIR
\end{quote}
\end{minipage} & 0.2528 & 0.0192 & 7.60\% & 0.0382 \\
\begin{minipage}[t]{\linewidth}\raggedright
\begin{quote}
Rank IC
\end{quote}
\end{minipage} & 0.0427 & 0.0024 & 5.56\% & 0.0045 \\
Rank ICIR & 0.2453 & 0.0210 & 8.55\% & 0.0419 \\
\bottomrule()
\end{longtable}

\begin{quote}
To further validate the robustness of the mined factors and rule out the
possibility of data-snooping, we summarize the daily IC and Rank IC
series over 966 trading days (from 2022-01-04 to 2025-12-26). As shown
in Table 6, the factors generated by both Claude-4.5-Sonnet and
DeepSeek-v3.2 exhibit highly significant predictive power. The
t-statistics far exceed standard significance thresholds (\emph{p
\textless{}} 0\emph{.}001), and the positive information coefficient
days remain above 60\%, confirming that the performance gains are
statistically solid and stable across the out-of-sample period.
\end{quote}

Table 6: Daily IC and Rank IC Statistics (966 Trading Days, 2022--2025)

\begin{longtable}[]{@{}
  >{\raggedright\arraybackslash}p{(\columnwidth - 16\tabcolsep) * \real{0.1111}}
  >{\raggedright\arraybackslash}p{(\columnwidth - 16\tabcolsep) * \real{0.1111}}
  >{\raggedright\arraybackslash}p{(\columnwidth - 16\tabcolsep) * \real{0.1111}}
  >{\raggedright\arraybackslash}p{(\columnwidth - 16\tabcolsep) * \real{0.1111}}
  >{\raggedright\arraybackslash}p{(\columnwidth - 16\tabcolsep) * \real{0.1111}}
  >{\raggedright\arraybackslash}p{(\columnwidth - 16\tabcolsep) * \real{0.1111}}
  >{\raggedright\arraybackslash}p{(\columnwidth - 16\tabcolsep) * \real{0.1111}}
  >{\raggedright\arraybackslash}p{(\columnwidth - 16\tabcolsep) * \real{0.1111}}
  >{\raggedright\arraybackslash}p{(\columnwidth - 16\tabcolsep) * \real{0.1111}}@{}}
\toprule()
\begin{minipage}[b]{\linewidth}\raggedright
\textbf{Factor Library}
\end{minipage} & \begin{minipage}[b]{\linewidth}\raggedright
\begin{quote}
\textbf{Metric}
\end{quote}
\end{minipage} & \begin{minipage}[b]{\linewidth}\raggedright
\textbf{Mean}
\end{minipage} & \begin{minipage}[b]{\linewidth}\raggedright
\textbf{Median}
\end{minipage} & \begin{minipage}[b]{\linewidth}\raggedright
\textbf{Std}
\end{minipage} & \begin{minipage}[b]{\linewidth}\raggedright
\textbf{Positive Days}
\end{minipage} & \begin{minipage}[b]{\linewidth}\raggedright
\textbf{95\% CI}
\end{minipage} & \begin{minipage}[b]{\linewidth}\raggedright
\textbf{t-stat}
\end{minipage} & \begin{minipage}[b]{\linewidth}\raggedright
\textbf{p-value}
\end{minipage} \\
\midrule()
\endhead
\begin{minipage}[t]{\linewidth}\raggedright
\begin{quote}
Claude
\end{quote}
\end{minipage} & \begin{minipage}[t]{\linewidth}\raggedright
\begin{quote}
IC
\end{quote}
\end{minipage} & 0.0426 & 0.0513 & 0.1833 & 60.04\% & {[}0.0311,
0.0542{]} & 7.23 & 4.95e-13 \\
\begin{minipage}[t]{\linewidth}\raggedright
\begin{quote}
Claude
\end{quote}
\end{minipage} & Rank IC & 0.0409 & 0.0438 & 0.1827 & 60.04\% &
{[}0.0293, 0.0524{]} & 6.95 & 3.68e-12 \\
DeepSeek-v3.2 & \begin{minipage}[t]{\linewidth}\raggedright
\begin{quote}
IC
\end{quote}
\end{minipage} & 0.0459 & 0.0448 & 0.1711 & 60.97\% & {[}0.0348,
0.0544{]} & 7.93 & 2.22e-15 \\
DeepSeek-v3.2 & Rank IC & 0.0418 & 0.0403 & 0.1694 & 60.97\% &
{[}0.0311, 0.0525{]} & 7.67 & 1.73e-14 \\
\bottomrule()
\end{longtable}

\begin{quote}
The results show that different seed combinations introduce only limited
fluctuations and do not change the overall conclusion. In particular,
the standard deviations of IC and Rank IC are 0.0021 and 0.0024,
respectively, suggesting that QuantaAlpha's performance is not driven by
a fortuitously selected seed set. This supports the robustness of the
proposed evolutionary factor-mining process under different
initializations.

Beyond initialization, QuantaAlpha also reduces uncontrolled LLM
randomness through constrained self-evolution. The framework combines
investment-hypothesis guidance, trajectory-pool memory, feedback-driven
iteration, and diversity-preserving mutation, which expands the search
space while preserving continuity and directional consistency in the
evolutionary path. As shown in Appendix E and Figure 6, performance
improves over early iterations and stabilizes as the factor pool
matures, indicating that the evolutionary process remains controllable
rather than drifting randomly.
\end{quote}

15

\begin{longtable}[]{@{}
  >{\raggedright\arraybackslash}p{(\columnwidth - 2\tabcolsep) * \real{0.5000}}
  >{\raggedright\arraybackslash}p{(\columnwidth - 2\tabcolsep) * \real{0.5000}}@{}}
\toprule()
\begin{minipage}[b]{\linewidth}\raggedright
\begin{quote}
\textbf{A.6}
\end{quote}
\end{minipage} & \begin{minipage}[b]{\linewidth}\raggedright
\begin{quote}
\textbf{Transaction Cost Sensitivity Analysis}
\end{quote}
\end{minipage} \\
\midrule()
\endhead
\bottomrule()
\end{longtable}

\begin{quote}
In our main experiments, we adopt a standard transaction-cost setting
with a buying cost of 0.05\% (open\_cost=0.0005) and a selling cost of
0.15\% (close\_cost=0.0015), corresponding to a 0.20\% round-trip cost.
We further evaluate robustness to trading frictions by scaling the costs
to 1.5\emph{×} (0.30\% round-trip) and 2.0\emph{×} (0.40\% round-trip).
As shown in the sensitivity analysis, QuantaAlpha maintains stable
annualized returns and Sharpe ratios even under doubled transaction
costs. The robustness mainly stems from the low-turnover portfolio
construction induced by the TopkDropoutStrategy, which replaces only
about 10\% of holdings at each rebalance. These results indicate that
QuantaAlpha remains effective under conservative cost assumptions and is
not driven by excessive turnover.

\textbf{B Algorithm Configuration}

This section details the evolution algorithm parameters, factor
constraints, and trading strategy configura-tion.

QuantaAlpha employs an evolutionary algorithm with mutation and
crossover operations, with Light-GBM used as the downstream model for
factor-based prediction. In the Planning Phase of the algorithm, we set
ten parallel exploration directions to broaden the initial coverage of
research space. For the experimental setup, beyond the original round,
the algorithm follows a fixed iterative process: the main experiment
consists of 5 total iterations, and each iteration comprises one
Mutation phase followed by one Crossover phase---alternating between
these two phases to iteratively refine high-quality hypotheses.
Additionally, we set a rule that each hypothesis attempts to generate 3
factor expressions, ensuring a focused yet sufficient exploration of
factor candidates under each research direction.

To prevent overfitting and ensure interpretability, factor
expressions---built using the operators listed in Table 7---are
restricted by the following constraints: symbol length \emph{$\le$}250
characters, base features \emph{$\le$}6, free arguments ratio
\emph{\textless{}} 50

Table 7: List of Supported Operators for Factor Construction
\end{quote}

\begin{longtable}[]{@{}
  >{\raggedright\arraybackslash}p{(\columnwidth - 4\tabcolsep) * \real{0.3333}}
  >{\raggedright\arraybackslash}p{(\columnwidth - 4\tabcolsep) * \real{0.3333}}
  >{\raggedright\arraybackslash}p{(\columnwidth - 4\tabcolsep) * \real{0.3333}}@{}}
\toprule()
\begin{minipage}[b]{\linewidth}\raggedright
\begin{quote}
\textbf{Category}
\end{quote}
\end{minipage} & \begin{minipage}[b]{\linewidth}\raggedright
\begin{quote}
\textbf{Operators}
\end{quote}
\end{minipage} & \begin{minipage}[b]{\linewidth}\raggedright
\begin{quote}
\textbf{Description}
\end{quote}
\end{minipage} \\
\midrule()
\endhead
\multirow{4}{*}{\begin{minipage}[t]{\linewidth}\raggedright
\begin{quote}
Time-Series
\end{quote}
\end{minipage}} & DELTA, DELAY, TS\_MEAN, TS\_STD, TS\_VAR, TS\_MAX,
TS\_MIN, & \begin{minipage}[t]{\linewidth}\raggedright
\begin{quote}
Rolling statistics computed along time axis
\end{quote}
\end{minipage} \\
& TS\_SUM, TS\_RANK, TS\_CORR, TS\_COVARIANCE, TS\_ARGMAX, &
\multirow{3}{*}{\begin{minipage}[t]{\linewidth}\raggedright
\begin{quote}
per instrument
\end{quote}
\end{minipage}} \\
& TS\_ARGMIN, TS\_SKEW, TS\_KURT, TS\_PCTCHANGE, TS\_ZSCORE, \\
& \begin{minipage}[t]{\linewidth}\raggedright
\begin{quote}
TS\_QUANTILE
\end{quote}
\end{minipage} \\
\multirow{2}{*}{Cross-Sectional} & \multirow{2}{*}{RANK, ZSCORE, SCALE,
MEAN, STD, MEDIAN, MAX, MIN, SKEW, KURT} &
\begin{minipage}[t]{\linewidth}\raggedright
\begin{quote}
Statistics computed across stocks per date-
\end{quote}
\end{minipage} \\
& & \begin{minipage}[t]{\linewidth}\raggedright
\begin{quote}
time
\end{quote}
\end{minipage} \\
\begin{minipage}[t]{\linewidth}\raggedright
\begin{quote}
Mathematical
\end{quote}
\end{minipage} & \begin{minipage}[t]{\linewidth}\raggedright
\begin{quote}
ABS, SIGN, LOG, EXP, SQRT, POW, INV
\end{quote}
\end{minipage} & \begin{minipage}[t]{\linewidth}\raggedright
\begin{quote}
Element-wise mathematical functions
\end{quote}
\end{minipage} \\
\multirow{2}{*}{\begin{minipage}[t]{\linewidth}\raggedright
\begin{quote}
Technical
\end{quote}
\end{minipage}} & SMA, EMA, WMA, MACD, RSI, BB\_UPPER, BB\_LOWER,
DECAYLINEAR, &
\multirow{2}{*}{\begin{minipage}[t]{\linewidth}\raggedright
\begin{quote}
Common technical indicators
\end{quote}
\end{minipage}} \\
& \begin{minipage}[t]{\linewidth}\raggedright
\begin{quote}
REGBETA, REGRESI
\end{quote}
\end{minipage} \\
\begin{minipage}[t]{\linewidth}\raggedright
\begin{quote}
Logical
\end{quote}
\end{minipage} & \begin{minipage}[t]{\linewidth}\raggedright
\begin{quote}
GT, LT, GE, LE, AND, OR, WHERE
\end{quote}
\end{minipage} & \begin{minipage}[t]{\linewidth}\raggedright
\begin{quote}
Comparison and conditional operators
\end{quote}
\end{minipage} \\
\begin{minipage}[t]{\linewidth}\raggedright
\begin{quote}
Auxiliary
\end{quote}
\end{minipage} & \begin{minipage}[t]{\linewidth}\raggedright
\begin{quote}
COUNT, SUMIF, FILTER, PROD, HIGHDAY, LOWDAY
\end{quote}
\end{minipage} & Helper functions for complex expressions \\
\bottomrule()
\end{longtable}

\begin{quote}
We employ a TopkDropout strategy for portfolio construction, as detailed
in Table 8. On each trading day, stocks are ranked according to their
predicted scores; the \emph{n}drop lowest-scoring holdings are
liquidated and replaced with the highest-ranked candidates to maintain a
constant portfolio size with equal weighting.

\textbf{C Case Study: Factor Evolution Trajectory}

This appendix presents a detailed case study of factor evolution in
QuantaAlpha. We trace the complete trajectory of a representative
factor---\emph{Institutional\_Momentum\_Score\_20D}---through the
crossover phase, demonstrating how the evolutionary framework
synthesizes complementary market hypotheses from parent trajectories.

QuantaAlpha's evolution process operates in three phases: (1)
\textbf{Original} phase where initial hypotheses are generated, (2)
\textbf{Mutation} phase where existing trajectories are perturbed to
explore diversified strategies,
\end{quote}

16

Table 8: Parameters for the TopkDropout Trading Strategy

\begin{longtable}[]{@{}
  >{\raggedright\arraybackslash}p{(\columnwidth - 4\tabcolsep) * \real{0.3333}}
  >{\raggedright\arraybackslash}p{(\columnwidth - 4\tabcolsep) * \real{0.3333}}
  >{\raggedright\arraybackslash}p{(\columnwidth - 4\tabcolsep) * \real{0.3333}}@{}}
\toprule()
\begin{minipage}[b]{\linewidth}\raggedright
\begin{quote}
\textbf{Parameter}
\end{quote}
\end{minipage} & \begin{minipage}[b]{\linewidth}\raggedright
\begin{quote}
\textbf{Value}
\end{quote}
\end{minipage} & \begin{minipage}[b]{\linewidth}\raggedright
\begin{quote}
\textbf{Description}
\end{quote}
\end{minipage} \\
\midrule()
\endhead
\bottomrule()
\end{longtable}

\emph{Portfolio}

\begin{longtable}[]{@{}
  >{\raggedright\arraybackslash}p{(\columnwidth - 4\tabcolsep) * \real{0.3333}}
  >{\raggedright\arraybackslash}p{(\columnwidth - 4\tabcolsep) * \real{0.3333}}
  >{\raggedright\arraybackslash}p{(\columnwidth - 4\tabcolsep) * \real{0.3333}}@{}}
\toprule()
\begin{minipage}[b]{\linewidth}\raggedright
\begin{quote}
topk
\end{quote}
\end{minipage} & \begin{minipage}[b]{\linewidth}\raggedright
50
\end{minipage} & \begin{minipage}[b]{\linewidth}\raggedright
Number of stocks held
\end{minipage} \\
\midrule()
\endhead
\begin{minipage}[t]{\linewidth}\raggedright
\begin{quote}
n\_drop
\end{quote}
\end{minipage} & 5 & Stocks dropped per rebalance \\
\bottomrule()
\end{longtable}

\emph{Transaction Costs}

\begin{longtable}[]{@{}
  >{\raggedright\arraybackslash}p{(\columnwidth - 4\tabcolsep) * \real{0.3333}}
  >{\raggedright\arraybackslash}p{(\columnwidth - 4\tabcolsep) * \real{0.3333}}
  >{\raggedright\arraybackslash}p{(\columnwidth - 4\tabcolsep) * \real{0.3333}}@{}}
\toprule()
\begin{minipage}[b]{\linewidth}\raggedright
\begin{quote}
Buying Fee
\end{quote}
\end{minipage} & \begin{minipage}[b]{\linewidth}\raggedright
\begin{quote}
0.05\%
\end{quote}
\end{minipage} & \begin{minipage}[b]{\linewidth}\raggedright
\begin{quote}
Commission
\end{quote}
\end{minipage} \\
\midrule()
\endhead
\begin{minipage}[t]{\linewidth}\raggedright
\begin{quote}
Selling Fee
\end{quote}
\end{minipage} & \begin{minipage}[t]{\linewidth}\raggedright
\begin{quote}
0.15\%
\end{quote}
\end{minipage} & \begin{minipage}[t]{\linewidth}\raggedright
\begin{quote}
Commission + stamp duty
\end{quote}
\end{minipage} \\
\bottomrule()
\end{longtable}

\emph{Execution}

\begin{longtable}[]{@{}
  >{\raggedright\arraybackslash}p{(\columnwidth - 4\tabcolsep) * \real{0.3333}}
  >{\raggedright\arraybackslash}p{(\columnwidth - 4\tabcolsep) * \real{0.3333}}
  >{\raggedright\arraybackslash}p{(\columnwidth - 4\tabcolsep) * \real{0.3333}}@{}}
\toprule()
\begin{minipage}[b]{\linewidth}\raggedright
\begin{quote}
Deal Price
\end{quote}
\end{minipage} & \begin{minipage}[b]{\linewidth}\raggedright
\begin{quote}
Open
\end{quote}
\end{minipage} & \begin{minipage}[b]{\linewidth}\raggedright
Next-day opening price
\end{minipage} \\
\midrule()
\endhead
Limit Threshold & \begin{minipage}[t]{\linewidth}\raggedright
\begin{quote}
9.5\%
\end{quote}
\end{minipage} & \begin{minipage}[t]{\linewidth}\raggedright
\begin{quote}
Price limit for halt
\end{quote}
\end{minipage} \\
\bottomrule()
\end{longtable}

\emph{Benchmark}

\begin{longtable}[]{@{}
  >{\raggedright\arraybackslash}p{(\columnwidth - 4\tabcolsep) * \real{0.3333}}
  >{\raggedright\arraybackslash}p{(\columnwidth - 4\tabcolsep) * \real{0.3333}}
  >{\raggedright\arraybackslash}p{(\columnwidth - 4\tabcolsep) * \real{0.3333}}@{}}
\toprule()
\begin{minipage}[b]{\linewidth}\raggedright
\begin{quote}
China
\end{quote}
\end{minipage} & \begin{minipage}[b]{\linewidth}\raggedright
\begin{quote}
SH000300/SH000905
\end{quote}
\end{minipage} & \begin{minipage}[b]{\linewidth}\raggedright
\begin{quote}
CSI 300/CSI500
\end{quote}
\end{minipage} \\
\midrule()
\endhead
\begin{minipage}[t]{\linewidth}\raggedright
\begin{quote}
U.S.
\end{quote}
\end{minipage} & \begin{minipage}[t]{\linewidth}\raggedright
\begin{quote}
SPX
\end{quote}
\end{minipage} & \begin{minipage}[t]{\linewidth}\raggedright
\begin{quote}
S\&P 500
\end{quote}
\end{minipage} \\
\bottomrule()
\end{longtable}

and (3) \textbf{Crossover} phase where high-performing parent
trajectories are combined to synthesize offspring with potentially
superior predictive power. The following factor card illustrates a
Crossover operation.

\begin{quote}
\textbf{C.1 Factor Identity}
\end{quote}

The factor card below presents the basic information of the evolved
factor, including its unique identifiers, evolution lineage, and
mathematical formulation.

\begin{longtable}[]{@{}
  >{\raggedright\arraybackslash}p{(\columnwidth - 2\tabcolsep) * \real{0.5000}}
  >{\raggedright\arraybackslash}p{(\columnwidth - 2\tabcolsep) * \real{0.5000}}@{}}
\toprule()
\multicolumn{2}{@{}>{\raggedright\arraybackslash}p{(\columnwidth - 2\tabcolsep) * \real{1.0000} + 2\tabcolsep}@{}}{%
\begin{minipage}[b]{\linewidth}\raggedright
\begin{quote}
\textbf{Institutional\_Momentum\_Score\_20D}
\end{quote}
\end{minipage}} \\
\midrule()
\endhead
\begin{minipage}[t]{\linewidth}\raggedright
\begin{quote}
\textbf{Factor ID:}
\end{quote}
\end{minipage} & c57cace576a95356 \\
\begin{minipage}[t]{\linewidth}\raggedright
\begin{quote}
\textbf{Trajectory ID:}
\end{quote}
\end{minipage} & df5a496878f4 \\
\textbf{Evolution Round:} & Round 10 \\
\textbf{Evolution Phase:} & \begin{minipage}[t]{\linewidth}\raggedright
\begin{minipage}[b]{\linewidth}\raggedright
\textbf{Crossover}
\end{minipage}
\end{minipage} \\
\begin{minipage}[t]{\linewidth}\raggedright
\begin{quote}
\textbf{Direction ID:}
\end{quote}
\end{minipage} & 6 \\
\bottomrule()
\end{longtable}

\begin{longtable}[]{@{}
  >{\raggedright\arraybackslash}p{(\columnwidth - 14\tabcolsep) * \real{0.1250}}
  >{\raggedright\arraybackslash}p{(\columnwidth - 14\tabcolsep) * \real{0.1250}}
  >{\raggedright\arraybackslash}p{(\columnwidth - 14\tabcolsep) * \real{0.1250}}
  >{\raggedright\arraybackslash}p{(\columnwidth - 14\tabcolsep) * \real{0.1250}}
  >{\raggedright\arraybackslash}p{(\columnwidth - 14\tabcolsep) * \real{0.1250}}
  >{\raggedright\arraybackslash}p{(\columnwidth - 14\tabcolsep) * \real{0.1250}}
  >{\raggedright\arraybackslash}p{(\columnwidth - 14\tabcolsep) * \real{0.1250}}
  >{\raggedright\arraybackslash}p{(\columnwidth - 14\tabcolsep) * \real{0.1250}}@{}}
\toprule()
\begin{minipage}[b]{\linewidth}\raggedright
\begin{quote}
\textbf{Factor Expression:}
\end{quote}
\end{minipage} & \begin{minipage}[b]{\linewidth}\raggedright
\end{minipage} & \begin{minipage}[b]{\linewidth}\raggedright
\end{minipage} & \begin{minipage}[b]{\linewidth}\raggedright
\end{minipage} & \begin{minipage}[b]{\linewidth}\raggedright
\end{minipage} & \begin{minipage}[b]{\linewidth}\raggedright
\end{minipage} & \begin{minipage}[b]{\linewidth}\raggedright
\end{minipage} & \begin{minipage}[b]{\linewidth}\raggedright
\end{minipage} \\
\midrule()
\endhead
\begin{minipage}[t]{\linewidth}\raggedright
\begin{quote}
RANK(TS\_CORR(DELTA(close,
\end{quote}
\end{minipage} & 1)/close, & DELTA(volume, & 1)/volume, & 20) & * &
TS\_MEAN((close & \begin{minipage}[t]{\linewidth}\raggedright
\begin{quote}
-
\end{quote}
\end{minipage} \\
\bottomrule()
\end{longtable}

\begin{longtable}[]{@{}
  >{\raggedright\arraybackslash}p{(\columnwidth - 0\tabcolsep) * \real{1.0000}}@{}}
\toprule()
\begin{minipage}[b]{\linewidth}\raggedright
\begin{quote}
open)/close, 5))
\end{quote}
\end{minipage} \\
\midrule()
\endhead
\bottomrule()
\end{longtable}

\begin{longtable}[]{@{}
  >{\raggedright\arraybackslash}p{(\columnwidth - 12\tabcolsep) * \real{0.1429}}
  >{\raggedright\arraybackslash}p{(\columnwidth - 12\tabcolsep) * \real{0.1429}}
  >{\raggedright\arraybackslash}p{(\columnwidth - 12\tabcolsep) * \real{0.1429}}
  >{\raggedright\arraybackslash}p{(\columnwidth - 12\tabcolsep) * \real{0.1429}}
  >{\raggedright\arraybackslash}p{(\columnwidth - 12\tabcolsep) * \real{0.1429}}
  >{\raggedright\arraybackslash}p{(\columnwidth - 12\tabcolsep) * \real{0.1429}}
  >{\raggedright\arraybackslash}p{(\columnwidth - 12\tabcolsep) * \real{0.1429}}@{}}
\toprule()
\multicolumn{7}{@{}>{\raggedright\arraybackslash}p{(\columnwidth - 12\tabcolsep) * \real{1.0000} + 12\tabcolsep}@{}}{%
\begin{minipage}[b]{\linewidth}\raggedright
\begin{quote}
\textbf{Mathematical Formulation:}
\end{quote}
\end{minipage}} \\
\midrule()
\endhead
\multirow{2}{*}{IMS20\emph{D} = RANK} & & & & & &
\multirow{2}{*}{\begin{minipage}[t]{\linewidth}\raggedright
\begin{quote}
\emph{,}
\end{quote}
\end{minipage}} \\
& \emph{$\rho$}20 & \begin{minipage}[t]{\linewidth}\raggedright
\begin{quote}
$\Delta$\emph{P}\\
\emph{P ,} $\Delta$\emph{V}
\end{quote}\strut
\end{minipage} & \emph{×} & \begin{minipage}[t]{\linewidth}\raggedright
\begin{quote}
\emph{C $-$O C}
\end{quote}
\end{minipage} & 5 \\
\multicolumn{7}{@{}>{\raggedright\arraybackslash}p{(\columnwidth - 12\tabcolsep) * \real{1.0000} + 12\tabcolsep}@{}}{%
where \emph{$\rho$}20(\emph{·, ·}) denotes the 20-day rolling correlation,
$\Delta$\emph{P/P} is the daily return, $\Delta$\emph{V/V} is the volume change ratio,
(\emph{·})5 is the 5-day moving average, and \emph{C} and \emph{O}
represent the closing and opening prices, respectively.} \\
\bottomrule()
\end{longtable}

\begin{longtable}[]{@{}
  >{\raggedright\arraybackslash}p{(\columnwidth - 0\tabcolsep) * \real{1.0000}}@{}}
\toprule()
\begin{minipage}[b]{\linewidth}\raggedright
\begin{quote}
\textbf{Factor Interpretation:}\\
This factor captures institutional-driven momentum by measuring two key
signals: (1) the correlation between price returns and volume changes,
which indicates coordinated institutional trading when positive; and (2)
the average intraday return pattern, reflecting institutional activity
that typically influences closing prices. The cross-sectional ranking
ensures comparability across stocks.
\end{quote}\strut
\end{minipage} \\
\midrule()
\endhead
\bottomrule()
\end{longtable}

17

\includegraphics[width=6.31944in,height=6.59722in]{vertopal_76ea1cc15b3d4460874f5e09e8e3ea79/media/image8.png}

\begin{longtable}[]{@{}
  >{\raggedright\arraybackslash}p{(\columnwidth - 2\tabcolsep) * \real{0.5000}}
  >{\raggedright\arraybackslash}p{(\columnwidth - 2\tabcolsep) * \real{0.5000}}@{}}
\toprule()
\begin{minipage}[b]{\linewidth}\raggedright
\begin{quote}
\textbf{C.2}
\end{quote}
\end{minipage} & \begin{minipage}[b]{\linewidth}\raggedright
\begin{quote}
\textbf{Evolution Lineage}
\end{quote}
\end{minipage} \\
\midrule()
\endhead
\bottomrule()
\end{longtable}

The crossover operation combines insights from two parent trajectories
with complementary market hypotheses. Parent 1 focuses on identifying
\emph{fragile momentum} driven by retail speculation, while Parent 2
targets \emph{sustainable momentum} supported by institutional activity.
The LLM synthesizes these complementary perspectives into a unified
framework.

\begin{quote}
\textbf{Evolution Information}

\textbf{Parent Trajectories:}

\textbf{Parent 1: 1e6d57e38e89}
\end{quote}

\begin{longtable}[]{@{}
  >{\raggedright\arraybackslash}p{(\columnwidth - 2\tabcolsep) * \real{0.5000}}
  >{\raggedright\arraybackslash}p{(\columnwidth - 2\tabcolsep) * \real{0.5000}}@{}}
\toprule()
\begin{minipage}[b]{\linewidth}\raggedright
\begin{quote}
\textbf{Round:}\\
\textbf{Phase:}\\
\textbf{Rank IC: IC:}\\
\textbf{IR:}
\end{quote}\strut
\end{minipage} & \begin{minipage}[b]{\linewidth}\raggedright
\begin{quote}
Round 9\\
Mutation\\
0.0216\\
0.0059\\
1.297
\end{quote}\strut
\end{minipage} \\
\midrule()
\endhead
\bottomrule()
\end{longtable}

\begin{quote}
\textbf{Core Hypothesis:}\\
When retail investors exhibit herd behavior and momentum chasing in
stocks with high social media activity, but accompanied by declining
institutional ownership and deteriorating fundamentals, the resulting
price momentum is unsustainable and leads to mean reversion.

\textbf{Parent 2: 47e0f0e55382}
\end{quote}

\begin{longtable}[]{@{}
  >{\raggedright\arraybackslash}p{(\columnwidth - 2\tabcolsep) * \real{0.5000}}
  >{\raggedright\arraybackslash}p{(\columnwidth - 2\tabcolsep) * \real{0.5000}}@{}}
\toprule()
\begin{minipage}[b]{\linewidth}\raggedright
\begin{quote}
\textbf{Round:}\\
\textbf{Phase:}\\
\textbf{Rank IC: IC:}\\
\textbf{IR:}
\end{quote}\strut
\end{minipage} & \begin{minipage}[b]{\linewidth}\raggedright
\begin{quote}
Round 8\\
Crossover\\
0.0246\\
0.0069\\
1.347
\end{quote}\strut
\end{minipage} \\
\midrule()
\endhead
\bottomrule()
\end{longtable}

\begin{quote}
\textbf{Core Hypothesis:}\\
A regime-adaptive structural momentum factor combining institutional
ownership-driven medium-term price trends with short-term microstructure
regime validation, where coordinated accumulation/distribution patterns
amplify momentum when confirmed by microstructure alignment.

\textbf{Evolution Path Diagram:}
\end{quote}

\begin{longtable}[]{@{}
  >{\raggedright\arraybackslash}p{(\columnwidth - 2\tabcolsep) * \real{0.5000}}
  >{\raggedright\arraybackslash}p{(\columnwidth - 2\tabcolsep) * \real{0.5000}}@{}}
\toprule()
\begin{minipage}[b]{\linewidth}\raggedright
\begin{quote}
Parent 1\\
Round 9 Mutation\\
Rank IC: 0.0216
\end{quote}\strut
\end{minipage} & \begin{minipage}[b]{\linewidth}\raggedright
\begin{quote}
Parent 2\\
Round 8 Crossover\\
Rank IC: 0.0246
\end{quote}\strut
\end{minipage} \\
\midrule()
\endhead
\bottomrule()
\end{longtable}

\begin{quote}
\textbf{Offspring Factor}\\
Round 8 Crossover\\
Rank IC: \textbf{0.0311}

\textbf{C.3 Synthesized Hypothesis}
\end{quote}

Through crossover, the LLM generates a new hypothesis that integrates
the complementary insights from both parents, rather than simply
averaging their factor expressions. This hypothesis-driven approach
ensures that the offspring factor captures genuinely novel market
dynamics.

\begin{longtable}[]{@{}
  >{\raggedright\arraybackslash}p{(\columnwidth - 0\tabcolsep) * \real{1.0000}}@{}}
\toprule()
\begin{minipage}[b]{\linewidth}\raggedright
\begin{quote}
\textbf{Hypothesis}
\end{quote}
\end{minipage} \\
\midrule()
\endhead
\bottomrule()
\end{longtable}

\begin{longtable}[]{@{}
  >{\raggedright\arraybackslash}p{(\columnwidth - 0\tabcolsep) * \real{1.0000}}@{}}
\toprule()
\begin{minipage}[b]{\linewidth}\raggedright
\begin{quote}
\textbf{Core Hypothesis:}\\
A regime-aware dual-source momentum factor that combines
institutional-driven structural momen-
\end{quote}\strut
\end{minipage} \\
\midrule()
\endhead
\bottomrule()
\end{longtable}

\begin{longtable}[]{@{}
  >{\raggedright\arraybackslash}p{(\columnwidth - 0\tabcolsep) * \real{1.0000}}@{}}
\toprule()
\begin{minipage}[b]{\linewidth}\raggedright
tum (validated by healthy microstructure) and retail-driven speculative
momentum (characterized
\end{minipage} \\
\midrule()
\endhead
\bottomrule()
\end{longtable}

18

\begin{longtable}[]{@{}
  >{\raggedright\arraybackslash}p{(\columnwidth - 0\tabcolsep) * \real{1.0000}}@{}}
\toprule()
\begin{minipage}[b]{\linewidth}\raggedright
by high attention and deteriorating fundamentals), dynamically weighted
by market volatility:
\end{minipage} \\
\midrule()
\endhead
\bottomrule()
\end{longtable}

amplifying institutional signals in stable regimes and retail reversal
signals in turbulent regimes,

\begin{longtable}[]{@{}
  >{\raggedright\arraybackslash}p{(\columnwidth - 2\tabcolsep) * \real{0.5000}}
  >{\raggedright\arraybackslash}p{(\columnwidth - 2\tabcolsep) * \real{0.5000}}@{}}
\toprule()
\multicolumn{2}{@{}>{\raggedright\arraybackslash}p{(\columnwidth - 2\tabcolsep) * \real{1.0000} + 2\tabcolsep}@{}}{%
\begin{minipage}[b]{\linewidth}\raggedright
will generate superior predictive returns.
\end{minipage}} \\
\midrule()
\endhead
\begin{minipage}[t]{\linewidth}\raggedright
\begin{quote}
\textbf{Component}
\end{quote}
\end{minipage} & \begin{minipage}[t]{\linewidth}\raggedright
\begin{quote}
\textbf{Description}
\end{quote}
\end{minipage} \\
\begin{minipage}[t]{\linewidth}\raggedright
\begin{quote}
\textbf{Observation}
\end{quote}
\end{minipage} & \begin{minipage}[t]{\linewidth}\raggedright
\begin{quote}
Parent strategies separately targeting institutional trends and retail
\end{quote}
\end{minipage} \\
\multicolumn{2}{@{}>{\raggedright\arraybackslash}p{(\columnwidth - 2\tabcolsep) * \real{1.0000} + 2\tabcolsep}@{}}{%
herding show moderate predictive power (Rank IC \emph{$\sim$}0.02--0.025),
suggesting combined signals could capture complementary market} \\
\bottomrule()
\end{longtable}

\begin{longtable}[]{@{}
  >{\raggedright\arraybackslash}p{(\columnwidth - 2\tabcolsep) * \real{0.5000}}
  >{\raggedright\arraybackslash}p{(\columnwidth - 2\tabcolsep) * \real{0.5000}}@{}}
\toprule()
\begin{minipage}[b]{\linewidth}\raggedright
\end{minipage} & \begin{minipage}[b]{\linewidth}\raggedright
\begin{quote}
dynamics.
\end{quote}
\end{minipage} \\
\midrule()
\endhead
\begin{minipage}[t]{\linewidth}\raggedright
\begin{quote}
\textbf{Justification}
\end{quote}
\end{minipage} & \begin{minipage}[t]{\linewidth}\raggedright
\begin{quote}
Sustainable price trends require institutional sponsorship and or-derly
trading, while retail-driven bubbles lack fundamental support
\end{quote}
\end{minipage} \\
\bottomrule()
\end{longtable}

and reverse under stress; a hybrid model exploiting both can en-

\begin{longtable}[]{@{}
  >{\raggedright\arraybackslash}p{(\columnwidth - 2\tabcolsep) * \real{0.5000}}
  >{\raggedright\arraybackslash}p{(\columnwidth - 2\tabcolsep) * \real{0.5000}}@{}}
\toprule()
\begin{minipage}[b]{\linewidth}\raggedright
\end{minipage} & \begin{minipage}[b]{\linewidth}\raggedright
\begin{quote}
hance robustness across market regimes.
\end{quote}
\end{minipage} \\
\midrule()
\endhead
\begin{minipage}[t]{\linewidth}\raggedright
\begin{quote}
\textbf{Domain Knowledge}
\end{quote}
\end{minipage} & \begin{minipage}[t]{\linewidth}\raggedright
\begin{quote}
Institutional accumulation with strong price-volume correlation and low
volatility indicates sustainable momentum; retail herding
\end{quote}
\end{minipage} \\
\bottomrule()
\end{longtable}

with declining institutional ownership and high volatility signals

\begin{longtable}[]{@{}
  >{\raggedright\arraybackslash}p{(\columnwidth - 4\tabcolsep) * \real{0.3333}}
  >{\raggedright\arraybackslash}p{(\columnwidth - 4\tabcolsep) * \real{0.3333}}
  >{\raggedright\arraybackslash}p{(\columnwidth - 4\tabcolsep) * \real{0.3333}}@{}}
\toprule()
\begin{minipage}[b]{\linewidth}\raggedright
\textbf{C.4}
\end{minipage} & \begin{minipage}[b]{\linewidth}\raggedright
\textbf{Backtest Performance}
\end{minipage} & \begin{minipage}[b]{\linewidth}\raggedright
\begin{quote}
fragile momentum prone to reversal.
\end{quote}
\end{minipage} \\
\midrule()
\endhead
\bottomrule()
\end{longtable}

After factor construction, QuantaAlpha automatically backtests the
generated factors using the Qlib framework. The results below compare
the offspring factor against both parent trajectories and the baseline,
demonstrating the effectiveness of the crossover operation.

\begin{longtable}[]{@{}
  >{\raggedright\arraybackslash}p{(\columnwidth - 6\tabcolsep) * \real{0.2500}}
  >{\raggedright\arraybackslash}p{(\columnwidth - 6\tabcolsep) * \real{0.2500}}
  >{\raggedright\arraybackslash}p{(\columnwidth - 6\tabcolsep) * \real{0.2500}}
  >{\raggedright\arraybackslash}p{(\columnwidth - 6\tabcolsep) * \real{0.2500}}@{}}
\toprule()
\begin{minipage}[b]{\linewidth}\raggedright
\begin{quote}
\textbf{Backtest Metrics}
\end{quote}
\end{minipage} & \begin{minipage}[b]{\linewidth}\raggedright
\begin{quote}
\textbf{Metric}
\end{quote}
\end{minipage} & \begin{minipage}[b]{\linewidth}\raggedright
\textbf{Offspring Factor}
\end{minipage} & \begin{minipage}[b]{\linewidth}\raggedright
\textbf{Baseline}
\end{minipage} \\
\midrule()
\endhead
\multirow{2}{*}{} & \begin{minipage}[t]{\linewidth}\raggedright
\begin{quote}
\textbf{IC}
\end{quote}
\end{minipage} & 0.0126 & 0.0058 \\
& \begin{minipage}[t]{\linewidth}\raggedright
\begin{quote}
\textbf{Rank IC}
\end{quote}
\end{minipage} & \textbf{0.0311} & 0.0220 \\
\multirow{3}{*}{} & \begin{minipage}[t]{\linewidth}\raggedright
\begin{quote}
\textbf{ARR (Excess)}
\end{quote}
\end{minipage} & 7.80\% & 5.20\% \\
& \begin{minipage}[t]{\linewidth}\raggedright
\begin{quote}
\textbf{IR}
\end{quote}
\end{minipage} & 0.963 & 0.973 \\
& \begin{minipage}[t]{\linewidth}\raggedright
\begin{quote}
\textbf{MDD (Excess)}
\end{quote}
\end{minipage} & \emph{$-$}11.37\% & \emph{$-$}7.30\% \\
\bottomrule()
\end{longtable}

\begin{longtable}[]{@{}
  >{\raggedright\arraybackslash}p{(\columnwidth - 8\tabcolsep) * \real{0.2000}}
  >{\raggedright\arraybackslash}p{(\columnwidth - 8\tabcolsep) * \real{0.2000}}
  >{\raggedright\arraybackslash}p{(\columnwidth - 8\tabcolsep) * \real{0.2000}}
  >{\raggedright\arraybackslash}p{(\columnwidth - 8\tabcolsep) * \real{0.2000}}
  >{\raggedright\arraybackslash}p{(\columnwidth - 8\tabcolsep) * \real{0.2000}}@{}}
\toprule()
\multicolumn{2}{@{}>{\raggedright\arraybackslash}p{(\columnwidth - 8\tabcolsep) * \real{0.4000} + 2\tabcolsep}}{%
\begin{minipage}[b]{\linewidth}\raggedright
\begin{quote}
\textbf{Detailed Statistics:}
\end{quote}
\end{minipage}} & \begin{minipage}[b]{\linewidth}\raggedright
\end{minipage} & \begin{minipage}[b]{\linewidth}\raggedright
\end{minipage} & \begin{minipage}[b]{\linewidth}\raggedright
\end{minipage} \\
\midrule()
\endhead
\multicolumn{2}{@{}>{\raggedright\arraybackslash}p{(\columnwidth - 8\tabcolsep) * \real{0.4000} + 2\tabcolsep}}{%
\textbf{Metric}} & \begin{minipage}[t]{\linewidth}\raggedright
\begin{quote}
\textbf{Value}
\end{quote}
\end{minipage} & \begin{minipage}[t]{\linewidth}\raggedright
\begin{quote}
\textbf{Metric}
\end{quote}
\end{minipage} & \begin{minipage}[t]{\linewidth}\raggedright
\begin{quote}
\textbf{Value}
\end{quote}
\end{minipage} \\
\multicolumn{2}{@{}>{\raggedright\arraybackslash}p{(\columnwidth - 8\tabcolsep) * \real{0.4000} + 2\tabcolsep}}{%
\multirow{3}{*}{\begin{minipage}[t]{\linewidth}\raggedright
\begin{quote}
\textbf{Daily Excess Return (w/o cost) Excess Return Std}\\
\textbf{L2 Train Loss}
\end{quote}\strut
\end{minipage}}} & 0.0328\% &
\multirow{3}{*}{\begin{minipage}[t]{\linewidth}\raggedright
\begin{quote}
\textbf{Daily Excess Return (w/ cost) Turnover (FFR)}\\
\textbf{L2 Valid Loss}
\end{quote}\strut
\end{minipage}} & \begin{minipage}[t]{\linewidth}\raggedright
\begin{quote}
0.0128\%
\end{quote}
\end{minipage} \\
& & \begin{minipage}[t]{\linewidth}\raggedright
\begin{quote}
0.52\%
\end{quote}
\end{minipage} & & \begin{minipage}[t]{\linewidth}\raggedright
\begin{quote}
100\%
\end{quote}
\end{minipage} \\
& & \begin{minipage}[t]{\linewidth}\raggedright
\begin{quote}
0.9936
\end{quote}
\end{minipage} & & 0.9962 \\
\textbf{C.5} & \begin{minipage}[t]{\linewidth}\raggedright
\begin{quote}
\textbf{Trajectory Summary}
\end{quote}
\end{minipage} & & & \\
\bottomrule()
\end{longtable}

After evaluating backtest results, the LLM provides structured summary.
This trajectory summary loop enables continuous improvement by learning
from both successes and failures.

\begin{longtable}[]{@{}
  >{\raggedright\arraybackslash}p{(\columnwidth - 0\tabcolsep) * \real{1.0000}}@{}}
\toprule()
\begin{minipage}[b]{\linewidth}\raggedright
\begin{quote}
\textbf{Evaluation \& Summary}
\end{quote}
\end{minipage} \\
\midrule()
\endhead
\bottomrule()
\end{longtable}

\begin{longtable}[]{@{}
  >{\raggedright\arraybackslash}p{(\columnwidth - 0\tabcolsep) * \real{1.0000}}@{}}
\toprule()
\begin{minipage}[b]{\linewidth}\raggedright
\begin{quote}
\textbf{Observations:}\\
The crossover operation demonstrates a trade-off between enhanced
predictive accuracy and
\end{quote}\strut
\end{minipage} \\
\midrule()
\endhead
\bottomrule()
\end{longtable}

\begin{longtable}[]{@{}
  >{\raggedright\arraybackslash}p{(\columnwidth - 0\tabcolsep) * \real{1.0000}}@{}}
\toprule()
\begin{minipage}[b]{\linewidth}\raggedright
\begin{quote}
increased risk exposure compared to the baseline:
\end{quote}
\end{minipage} \\
\midrule()
\endhead
\bottomrule()
\end{longtable}

\begin{longtable}[]{@{}
  >{\raggedright\arraybackslash}p{(\columnwidth - 0\tabcolsep) * \real{1.0000}}@{}}
\toprule()
\begin{minipage}[b]{\linewidth}\raggedright
\begin{quote}
•Significant improvement in annualized excess return and predictive
metrics (IC and Rank
\end{quote}
\end{minipage} \\
\midrule()
\endhead
\bottomrule()
\end{longtable}

19

\begin{longtable}[]{@{}
  >{\raggedright\arraybackslash}p{(\columnwidth - 0\tabcolsep) * \real{1.0000}}@{}}
\toprule()
\begin{minipage}[b]{\linewidth}\raggedright
\begin{quote}
IC), validating the effectiveness of synthesizing dual-source momentum
signals.
\end{quote}
\end{minipage} \\
\midrule()
\endhead
\bottomrule()
\end{longtable}

\begin{quote}
•Increased maximum drawdown and a marginal decline in the Information
Ratio, suggesting
\end{quote}

\begin{longtable}[]{@{}
  >{\raggedright\arraybackslash}p{(\columnwidth - 0\tabcolsep) * \real{1.0000}}@{}}
\toprule()
\begin{minipage}[b]{\linewidth}\raggedright
\begin{quote}
that the offspring factor introduces higher volatility during certain
market regimes.
\end{quote}
\end{minipage} \\
\midrule()
\endhead
\bottomrule()
\end{longtable}

\begin{longtable}[]{@{}
  >{\raggedright\arraybackslash}p{(\columnwidth - 0\tabcolsep) * \real{1.0000}}@{}}
\toprule()
\begin{minipage}[b]{\linewidth}\raggedright
\begin{quote}
\textbf{Hypothesis Evaluation:}
\end{quote}
\end{minipage} \\
\midrule()
\endhead
\bottomrule()
\end{longtable}

\begin{longtable}[]{@{}
  >{\raggedright\arraybackslash}p{(\columnwidth - 0\tabcolsep) * \real{1.0000}}@{}}
\toprule()
\begin{minipage}[b]{\linewidth}\raggedright
\begin{quote}
Results partially support the hypothesis. Improved annualized return and
IC suggest that combining institutional and retail momentum signals has
merit. However, deterioration in risk metrics indicates that without
proper regime-adaptive weighting, the combined signals may amplify risks
during turbulent periods. The full hypothesis requires all three
components (institutional momentum, retail herding reversal,
volatility-adaptive weighting) to work effectively.
\end{quote}
\end{minipage} \\
\midrule()
\endhead
\bottomrule()
\end{longtable}

\begin{longtable}[]{@{}
  >{\raggedright\arraybackslash}p{(\columnwidth - 0\tabcolsep) * \real{1.0000}}@{}}
\toprule()
\begin{minipage}[b]{\linewidth}\raggedright
\begin{quote}
\textbf{Decision: REJECTED} for direct deployment.
\end{quote}
\end{minipage} \\
\midrule()
\endhead
\bottomrule()
\end{longtable}

\begin{longtable}[]{@{}
  >{\raggedright\arraybackslash}p{(\columnwidth - 0\tabcolsep) * \real{1.0000}}@{}}
\toprule()
\begin{minipage}[b]{\linewidth}\raggedright
\begin{quote}
\textbf{Recommendations:}
\end{quote}
\end{minipage} \\
\midrule()
\endhead
\bottomrule()
\end{longtable}

\begin{quote}
1. Use 20-day price-volume correlation as institutional momentum proxy;

2. Use 5-day average intraday returns as retail attention proxy;

3. Add volatility regime indicator (recent/historical volatility ratio)
for dynamic weighting.
\end{quote}

\begin{longtable}[]{@{}
  >{\raggedright\arraybackslash}p{(\columnwidth - 2\tabcolsep) * \real{0.5000}}
  >{\raggedright\arraybackslash}p{(\columnwidth - 2\tabcolsep) * \real{0.5000}}@{}}
\toprule()
\begin{minipage}[b]{\linewidth}\raggedright
\end{minipage} & \begin{minipage}[b]{\linewidth}\raggedright
\begin{quote}
This summary will inform the next mutation round, guiding the LLM to
simplify the factor expression while preserving the core dual-source
concept.
\end{quote}
\end{minipage} \\
\midrule()
\endhead
\textbf{D} & \begin{minipage}[t]{\linewidth}\raggedright
\begin{quote}
\textbf{Market Style, Factor Semantics, and Alpha Decay}
\end{quote}
\end{minipage} \\
\bottomrule()
\end{longtable}

This appendix provides a factor-level diagnosis of the alpha decay
observed in 2023, grounding the analysis jointly in \emph{market style
changes}, \emph{factor semantics}, and \emph{empirical factor
statistics}. Rather than enumerating factor formulas, we focus on which
information channels remain predictive under the 2022--2023 regime
transition and why.

\begin{quote}
\textbf{D.1 Market style transition and its implications}
\end{quote}

The CSI 300 market undergoes a pronounced style transition in 2023
relative to the training period (2016--2020). The earlier regime is
dominated by large-cap ``core assets,'' with high institutional
partici-pation, smooth intraday dynamics, and persistent short-horizon
trends. Under such conditions, classical momentum, mean-reversion, and
exhaustion-based factors are effective.

In 2023, leadership rotates toward small-cap and thematic stocks. This
shift is accompanied by increased intraday noise, frequent overnight
gaps driven by call auctions, and rapid cross-style liquidity
re-allocation. These changes weaken trend persistence and amplify
path-dependent noise, directly challenging factor constructions that
rely on stable intraday structure or fast mean reversion.

\begin{quote}
\textbf{D.2 Factor semantics aligned with market microstructure}
\end{quote}

The empirical contrast between QuantaAlpha (QA) and AlphaAgent (AA) in
2023 reflects differences in the \emph{semantic composition} of their
factor libraries.

\textbf{Overnight and auction information.} Gap-based factors aggregate
information released during non-trading hours and price discovery in the
opening auction. As shown in Table 9, multiple QA factors from this
channel occupy the right tail of the 2023 Rank IC distribution, while
maintaining near-complete coverage (Table 10). This indicates that
overnight information becomes a dominant and stable signal when intraday
predictability deteriorates.

\begin{quote}
\textbf{Volatility structure and range-based signals.} Range deviation
factors conditioned on volatility cluster-ing capture abnormal
variability rather than directional trends. In 2023, these signals
remain predictive despite elevated noise, consistent with the
persistence of volatility clustering across market styles. Their
contribution is reflected in the heavier positive Rank IC tail of QA
relative to AA.

\textbf{Trend quality and liquidity re-rating.} QA emphasizes trend
continuation only when supported by low residual volatility and
improving liquidity (e.g., rising dollar volume with limited price
impact). This
\end{quote}

20

\begin{quote}
Table 9: Representative factors and their performance metrics on the CSI
300 index in 2023.
\end{quote}

\begin{longtable}[]{@{}
  >{\raggedright\arraybackslash}p{(\columnwidth - 6\tabcolsep) * \real{0.2500}}
  >{\raggedright\arraybackslash}p{(\columnwidth - 6\tabcolsep) * \real{0.2500}}
  >{\raggedright\arraybackslash}p{(\columnwidth - 6\tabcolsep) * \real{0.2500}}
  >{\raggedright\arraybackslash}p{(\columnwidth - 6\tabcolsep) * \real{0.2500}}@{}}
\toprule()
\begin{minipage}[b]{\linewidth}\raggedright
\begin{quote}
\textbf{Factor (representative)}
\end{quote}
\end{minipage} & \begin{minipage}[b]{\linewidth}\raggedright
\textbf{Rank IC}
\end{minipage} & \begin{minipage}[b]{\linewidth}\raggedright
\textbf{IC}
\end{minipage} & \begin{minipage}[b]{\linewidth}\raggedright
\begin{quote}
\textbf{Interpretation (short)}
\end{quote}
\end{minipage} \\
\midrule()
\endhead
\bottomrule()
\end{longtable}

\emph{QA --- strong performers (overnight/auction, trend-quality, and
liquidity re-rating)}

\begin{longtable}[]{@{}
  >{\raggedright\arraybackslash}p{(\columnwidth - 6\tabcolsep) * \real{0.2500}}
  >{\raggedright\arraybackslash}p{(\columnwidth - 6\tabcolsep) * \real{0.2500}}
  >{\raggedright\arraybackslash}p{(\columnwidth - 6\tabcolsep) * \real{0.2500}}
  >{\raggedright\arraybackslash}p{(\columnwidth - 6\tabcolsep) * \real{0.2500}}@{}}
\toprule()
\begin{minipage}[b]{\linewidth}\raggedright
\begin{quote}
GapZ10\_Overnight\_vs\_TR
\end{quote}
\end{minipage} & \begin{minipage}[b]{\linewidth}\raggedright
0.0793
\end{minipage} & \begin{minipage}[b]{\linewidth}\raggedright
0.0335
\end{minipage} & \begin{minipage}[b]{\linewidth}\raggedright
Normalized overnight gap magnitude relative to recent true range;
captures auction-
\end{minipage} \\
\midrule()
\endhead
\multirow{2}{*}{\begin{minipage}[t]{\linewidth}\raggedright
\begin{quote}
Gap\_IntradayAcceptanceScore\_20D
\end{quote}
\end{minipage}} & \multirow{3}{*}{0.0744} & \multirow{3}{*}{0.0330} &
\begin{minipage}[t]{\linewidth}\raggedright
\begin{quote}
driven shocks and subsequent adjustment.
\end{quote}
\end{minipage} \\
& & & `Acceptance vs. rejection'' of an overnight gap using intraday
direction, scaled by \\
\multirow{3}{*}{\begin{minipage}[t]{\linewidth}\raggedright
\begin{quote}
Gap\_IntradayAcceptance\_VolWeighted\_20D
\end{quote}
\end{minipage}} & & &
\multirow{2}{*}{\begin{minipage}[t]{\linewidth}\raggedright
\begin{quote}
recent volatility.
\end{quote}
\end{minipage}} \\
& \multirow{2}{*}{0.0606} & \multirow{2}{*}{0.0314} \\
& & & Gap acceptance score weighted by abnormal volume; emphasizes
information-rich \\
\multirow{2}{*}{\begin{minipage}[t]{\linewidth}\raggedright
\begin{quote}
CleanTrend\_Continuation\_Score\_RS10\_WVMA5
\end{quote}
\end{minipage}} & \multirow{3}{*}{0.0590} & \multirow{3}{*}{0.0267} &
\begin{minipage}[t]{\linewidth}\raggedright
\begin{quote}
openings with high participation.
\end{quote}
\end{minipage} \\
& & & Trend continuation conditioned on high trend quality (low residual
noise) and muted \\
\multirow{4}{*}{\begin{minipage}[t]{\linewidth}\raggedright
\begin{quote}
OrderlyTrend\_x\_Absorption\_10D\_5D\_20D
\end{quote}
\end{minipage}} & & &
\multirow{2}{*}{\begin{minipage}[t]{\linewidth}\raggedright
\begin{quote}
intraday/volume pressure.
\end{quote}
\end{minipage}} \\
& \multirow{3}{*}{0.0465} & \multirow{3}{*}{0.0271} \\
& & & `Orderly'' short-horizon trend cross-validated by liquidity
absorption (high dollar \\
& & & \begin{minipage}[t]{\linewidth}\raggedright
\begin{quote}
volume with low price impact).
\end{quote}
\end{minipage} \\
\bottomrule()
\end{longtable}

\emph{QA --- weak performers (overly rigid gates or noisy-path proxies
under rotation)}

\begin{longtable}[]{@{}
  >{\raggedright\arraybackslash}p{(\columnwidth - 6\tabcolsep) * \real{0.2500}}
  >{\raggedright\arraybackslash}p{(\columnwidth - 6\tabcolsep) * \real{0.2500}}
  >{\raggedright\arraybackslash}p{(\columnwidth - 6\tabcolsep) * \real{0.2500}}
  >{\raggedright\arraybackslash}p{(\columnwidth - 6\tabcolsep) * \real{0.2500}}@{}}
\toprule()
\begin{minipage}[b]{\linewidth}\raggedright
\begin{quote}
KineticLength\_AbsRetSum\_Z\_10D
\end{quote}
\end{minipage} & \begin{minipage}[b]{\linewidth}\raggedright
-0.0720
\end{minipage} & \begin{minipage}[b]{\linewidth}\raggedright
-0.0246
\end{minipage} & \begin{minipage}[b]{\linewidth}\raggedright
Path-length proxy (choppiness): can behave like a noise detector, but
may invert
\end{minipage} \\
\midrule()
\endhead
\multirow{2}{*}{\begin{minipage}[t]{\linewidth}\raggedright
\begin{quote}
Drawdown\_Gated\_NegCorr\_60D\_20D\_thr20pct
\end{quote}
\end{minipage}} & \multirow{3}{*}{-0.0282} & \multirow{3}{*}{-0.0095} &
\begin{minipage}[t]{\linewidth}\raggedright
\begin{quote}
under fast style rotation.
\end{quote}
\end{minipage} \\
& & & Hard regime gate based on deep drawdown; brittle when drawdowns
cluster and \\
\multirow{4}{*}{\begin{minipage}[t]{\linewidth}\raggedright
\begin{quote}
HighClose\_Shock\_With\_VolSync\_60\_20
\end{quote}
\end{minipage}} & & &
\multirow{2}{*}{\begin{minipage}[t]{\linewidth}\raggedright
\begin{quote}
cross-sectional regimes shift quickly.
\end{quote}
\end{minipage}} \\
& \multirow{3}{*}{-0.0274} & \multirow{3}{*}{-0.0090} \\
& & & `Shock-day'' breakout quality (close-in-range, range shock,
return--volume sync); \\
& & & \begin{minipage}[t]{\linewidth}\raggedright
\begin{quote}
sensitive to regime-dependent follow-through.
\end{quote}
\end{minipage} \\
\bottomrule()
\end{longtable}

\emph{AA --- strong performers (exhaustion/climax-style reversals)}

\begin{longtable}[]{@{}
  >{\raggedright\arraybackslash}p{(\columnwidth - 6\tabcolsep) * \real{0.2500}}
  >{\raggedright\arraybackslash}p{(\columnwidth - 6\tabcolsep) * \real{0.2500}}
  >{\raggedright\arraybackslash}p{(\columnwidth - 6\tabcolsep) * \real{0.2500}}
  >{\raggedright\arraybackslash}p{(\columnwidth - 6\tabcolsep) * \real{0.2500}}@{}}
\toprule()
\begin{minipage}[b]{\linewidth}\raggedright
\begin{quote}
Exhaustion\_Intensity\_Index\_10D
\end{quote}
\end{minipage} &
\multirow{2}{*}{\begin{minipage}[b]{\linewidth}\raggedright
0.0323
\end{minipage}} &
\multirow{2}{*}{\begin{minipage}[b]{\linewidth}\raggedright
0.0159
\end{minipage}} & \begin{minipage}[b]{\linewidth}\raggedright
Price displacement over 60D interacted with volume intensity ratio;
targets potential
\end{minipage} \\
\multirow{3}{*}{\begin{minipage}[b]{\linewidth}\raggedright
\begin{quote}
Climax\_Exhaustion\_Intensity
\end{quote}
\end{minipage}} & & &
\multirow{2}{*}{\begin{minipage}[b]{\linewidth}\raggedright
\begin{quote}
exhaustion and reversal.
\end{quote}
\end{minipage}} \\
& \multirow{3}{*}{\begin{minipage}[b]{\linewidth}\raggedright
0.0242
\end{minipage}} &
\multirow{3}{*}{\begin{minipage}[b]{\linewidth}\raggedright
0.0160
\end{minipage}} \\
& & & \begin{minipage}[b]{\linewidth}\raggedright
Variant using short-horizon volume climax vs. long-horizon baseline;
aims to identify
\end{minipage} \\
\multirow{4}{*}{\begin{minipage}[b]{\linewidth}\raggedright
\begin{quote}
Exhaustion\_Volume\_Instability\_Index
\end{quote}
\end{minipage}} & & &
\multirow{2}{*}{\begin{minipage}[b]{\linewidth}\raggedright
\begin{quote}
capitulation-like turns.
\end{quote}
\end{minipage}} \\
& \multirow{3}{*}{\begin{minipage}[b]{\linewidth}\raggedright
0.0121
\end{minipage}} &
\multirow{3}{*}{\begin{minipage}[b]{\linewidth}\raggedright
0.0117
\end{minipage}} \\
& & & \begin{minipage}[b]{\linewidth}\raggedright
Trend deviation combined with volume instability; highlights fragile
price levels
\end{minipage} \\
& & & \begin{minipage}[b]{\linewidth}\raggedright
\begin{quote}
supported by unstable liquidity.
\end{quote}
\end{minipage} \\
\midrule()
\endhead
\bottomrule()
\end{longtable}

\emph{AA --- weak performers (bottom-fishing under non-stationary
liquidity)}

\begin{longtable}[]{@{}
  >{\raggedright\arraybackslash}p{(\columnwidth - 6\tabcolsep) * \real{0.2500}}
  >{\raggedright\arraybackslash}p{(\columnwidth - 6\tabcolsep) * \real{0.2500}}
  >{\raggedright\arraybackslash}p{(\columnwidth - 6\tabcolsep) * \real{0.2500}}
  >{\raggedright\arraybackslash}p{(\columnwidth - 6\tabcolsep) * \real{0.2500}}@{}}
\toprule()
\begin{minipage}[b]{\linewidth}\raggedright
\begin{quote}
Relative\_Volume\_Calm\_Reversal
\end{quote}
\end{minipage} &
\multirow{2}{*}{\begin{minipage}[b]{\linewidth}\raggedright
-0.0279
\end{minipage}} &
\multirow{2}{*}{\begin{minipage}[b]{\linewidth}\raggedright
-0.0188
\end{minipage}} & \begin{minipage}[b]{\linewidth}\raggedright
`Quiet-volume'' regimes multiplied by momentum divergence; may fail when
liquid-
\end{minipage} \\
\multirow{3}{*}{\begin{minipage}[b]{\linewidth}\raggedright
\begin{quote}
Volume\_Stability\_Momentum\_Divergence\_40D
\end{quote}
\end{minipage}} & & &
\multirow{2}{*}{\begin{minipage}[b]{\linewidth}\raggedright
\begin{quote}
ity conditions change abruptly.
\end{quote}
\end{minipage}} \\
& \multirow{2}{*}{\begin{minipage}[b]{\linewidth}\raggedright
-0.0247
\end{minipage}} &
\multirow{2}{*}{\begin{minipage}[b]{\linewidth}\raggedright
-0.0155
\end{minipage}} \\
& & & \begin{minipage}[b]{\linewidth}\raggedright
Robust volume-stability proxy (MAD) times momentum spread; sensitive to
turnover
\end{minipage} \\
\midrule()
\endhead
\multirow{3}{*}{\begin{minipage}[t]{\linewidth}\raggedright
\begin{quote}
LVR\_Bottom\_Fishing\_20D
\end{quote}
\end{minipage}} & \multirow{3}{*}{-0.0190} & \multirow{3}{*}{-0.0144} &
\begin{minipage}[t]{\linewidth}\raggedright
\begin{quote}
regime changes.
\end{quote}
\end{minipage} \\
& & & `Bottom-fishing'' reversal with intraday rejection and volume
surge; vulnerable when \\
& & & \begin{minipage}[t]{\linewidth}\raggedright
\begin{quote}
reversals are short-lived and crowded.
\end{quote}
\end{minipage} \\
\bottomrule()
\end{longtable}

\begin{quote}
Table 10: Summary statistics of annual factor predictability on the CSI
300 index in 2023.
\end{quote}

\begin{longtable}[]{@{}
  >{\raggedright\arraybackslash}p{(\columnwidth - 4\tabcolsep) * \real{0.3333}}
  >{\raggedright\arraybackslash}p{(\columnwidth - 4\tabcolsep) * \real{0.3333}}
  >{\raggedright\arraybackslash}p{(\columnwidth - 4\tabcolsep) * \real{0.3333}}@{}}
\toprule()
\begin{minipage}[b]{\linewidth}\raggedright
\begin{quote}
\textbf{Metric}
\end{quote}
\end{minipage} & \begin{minipage}[b]{\linewidth}\raggedright
\textbf{QA}
\end{minipage} & \begin{minipage}[b]{\linewidth}\raggedright
\textbf{AA}
\end{minipage} \\
\midrule()
\endhead
Coverage ratio (valid metrics) & 0.98 & 0.80 \\
\begin{minipage}[t]{\linewidth}\raggedright
\begin{quote}
Share with Rank IC \emph{\textgreater{}} 0
\end{quote}
\end{minipage} & 0.626 & 0.594 \\
\begin{minipage}[t]{\linewidth}\raggedright
\begin{quote}
Mean Rank IC
\end{quote}
\end{minipage} & 0.0057 & 0.0012 \\
\begin{minipage}[t]{\linewidth}\raggedright
\begin{quote}
Max Rank IC
\end{quote}
\end{minipage} & 0.0793 & 0.0323 \\
\begin{minipage}[t]{\linewidth}\raggedright
\begin{quote}
Min Rank IC
\end{quote}
\end{minipage} & -0.0720 & -0.0279 \\
\begin{minipage}[t]{\linewidth}\raggedright
\begin{quote}
Share with Rank IC \emph{\textgreater{}} 0\emph{.}03
\end{quote}
\end{minipage} & 0.102 & 0.0156 \\
\begin{minipage}[t]{\linewidth}\raggedright
\begin{quote}
Share with Rank IC \emph{\textgreater{}} 0\emph{.}05
\end{quote}
\end{minipage} & 0.0272 & 0.0000 \\
\begin{minipage}[t]{\linewidth}\raggedright
\begin{quote}
Mean IC
\end{quote}
\end{minipage} & 0.0044 & 0.0015 \\
\bottomrule()
\end{longtable}

\begin{quote}
conditioning filters out noise-driven pseudo-trends prevalent in
small-cap rotation, explaining why QA retains a higher fraction of
strong-performing factors, whereas raw momentum proxies degrade.

\textbf{Reversal and exhaustion signals.} AA and baseline methods rely
primarily on volume--price exhaustion and bottom-fishing logic. While
some such factors remain marginally positive, Table 10 shows weaker Rank
IC tails and lower coverage, indicating that reversal hypotheses become
unstable under non-stationary liquidity and thematic crowding.

\textbf{D.3 Factor statistics as supporting evidence}

The factor-level statistics in 2023 provide quantitative support for the
semantic interpretation above. QA achieves substantially higher coverage
and a heavier right tail of Rank IC, including a non-trivial share of
factors exceeding moderate predictability thresholds, whereas AA
exhibits compressed tails and fewer robust signals. These differences
are consistent with the dominance of overnight, volatility-structure,
and trend-quality channels under the 2023 market style.
\end{quote}

21

\begin{longtable}[]{@{}
  >{\raggedright\arraybackslash}p{(\columnwidth - 2\tabcolsep) * \real{0.5000}}
  >{\raggedright\arraybackslash}p{(\columnwidth - 2\tabcolsep) * \real{0.5000}}@{}}
\toprule()
\begin{minipage}[b]{\linewidth}\raggedright
\begin{quote}
\textbf{D.4}
\end{quote}
\end{minipage} & \begin{minipage}[b]{\linewidth}\raggedright
\begin{quote}
\textbf{Role of factor diversity and mutation}
\end{quote}
\end{minipage} \\
\midrule()
\endhead
\bottomrule()
\end{longtable}

\begin{quote}
A key distinction is that QA explicitly encourages semantic diversity
through its factor mutation mecha-nism. By generating and recombining
heterogeneous primitives across multiple information channels, QA avoids
concentration on a single market hypothesis. This diversity increases
the probability that a subset of factors remains aligned with the
prevailing market microstructure after a style transition, thereby
mitigating regime-specific alpha decay.

\textbf{D.5 Summary}

Overall, the 2023 diagnosis indicates that alpha robustness under
distribution shift is driven by alignment between market style and
factor semantics, supported by sufficient diversity across information
channels. QA's ability to maintain performance arises not from fitting a
specific regime, but from sustaining a broad and adaptive factor
population whose predictive subsets persist across market style
transitions.

\textbf{D.6 Stress-Period Analysis (High-Volatility Episode in April
2025)}

To evaluate the robustness of our generated factor pool under extreme
market conditions, we analyze the strategy's performance during the
severe market stress period in April 2025. This episode was
characterized by sharp index drawdowns, elevated cross-sectional
volatility, and rapid style rotations, which challenged the stability of
traditional quantitative factors and baselines.

During this black-swan regime, traditional factors often suffer from
crowding and severe capacity con-straints. In contrast, thanks to the
diversity-preserving mutation and crossover mechanisms, QuantaAlpha
maintains a more orthogonal and robust factor pool. Consequently, the
strategy exhibited a significantly lower maximum drawdown and faster
recovery during the April 2025 volatility spike compared to Al-phaAgent
and RD-Agent. This stress-period analysis confirms that evolving
complete research trajectories better equips the model to navigate
structural market shifts and extreme tail risks.

\textbf{E Iteration Study: Convergence of Factor-Pool Performance
(CSI300)}
\end{quote}

We run QuantaAlpha for a total of 15 iterations on CSI300, using
Deepseek-V3.2 as the backbone LLM for all agents. To improve mining
throughput, we set the factor-complexity constraints to symbol length

\emph{Crossover} phase.At the end of each iteration, we maintain a
global high-quality factor pool using all\emph{$\le$}200 and base features
\emph{$\le$}4. Each iteration consists of two phases: a \emph{Mutation}
phase followed by a

\begin{quote}
factors generated up to that iteration. The maintenance rule is greedy
and RankIC-driven. The RankIC gate for pool admission is computed
exclusively on the 2021 validation period; the 2022--2025 test window is
never accessed during evolution or pool construction. We sort all
candidate factors by RankIC (descending) and add them into the pool in
order, subject to a redundancy constraint. A factor is admitted only if
its absolute correlation with every factor already in the pool is below
0.7. The pool size is capped at 50\% of the total number of mined
factors up to the current iteration.

Figure 6 illustrates how QuantaAlpha refines hypotheses and factor
expressions across iterations via the \emph{mutation} and
\emph{crossover} phases. Each panel corresponds to one step along an
evolution trace and summarizes: (i) the iteration and phase, (ii) the
current hypothesis, (iii) a representative key factor generated under
that hypothesis (with its expression), and (iv) the LLM-based
feedback/observation that guides the next update. The highlighted
factors in the figure are therefore representative nodes on different
traces across iterations. We observe that the predictive performance
does not increase linearly with the number of iterations. The strategy's
return capacity and risk control ability are gradually optimized,
reaching a balanced high level between return and maximum drawdown
around the 11th to 12th iterations. Subsequent additional iterations
fail to bring significant performance improvement; instead, they may
introduce redundant information through newly generated factors, leading
to reduced strategy robustness and deteriorated drawdown performance.
Overall, the 11th to 12th iterations (approximately 350 factors in
total) represent the optimal trade-off point between return level and
risk control effect under this experimental setup.
\end{quote}

22

\includegraphics[width=6.29861in,height=3.45296in]{vertopal_76ea1cc15b3d4460874f5e09e8e3ea79/media/image9.png}\includegraphics[width=6.26389in,height=3.43056in]{vertopal_76ea1cc15b3d4460874f5e09e8e3ea79/media/image10.png}

\begin{longtable}[]{@{}
  >{\raggedright\arraybackslash}p{(\columnwidth - 10\tabcolsep) * \real{0.1667}}
  >{\raggedright\arraybackslash}p{(\columnwidth - 10\tabcolsep) * \real{0.1667}}
  >{\raggedright\arraybackslash}p{(\columnwidth - 10\tabcolsep) * \real{0.1667}}
  >{\raggedright\arraybackslash}p{(\columnwidth - 10\tabcolsep) * \real{0.1667}}
  >{\raggedright\arraybackslash}p{(\columnwidth - 10\tabcolsep) * \real{0.1667}}
  >{\raggedright\arraybackslash}p{(\columnwidth - 10\tabcolsep) * \real{0.1667}}@{}}
\toprule()
\begin{minipage}[b]{\linewidth}\raggedright
\textbf{Iter 2} \textbar{} \textbf{Mutation}
\end{minipage} & \begin{minipage}[b]{\linewidth}\raggedright
\begin{quote}
\textbf{Iter 2} \textbar{} \textbf{Crossover}
\end{quote}
\end{minipage} &
\multirow{15}{*}{\begin{minipage}[b]{\linewidth}\raggedright
\textbf{Factor Pool}
\end{minipage}} &
\multirow{15}{*}{\begin{minipage}[b]{\linewidth}\raggedright
27.85\%
\end{minipage}} &
\multirow{15}{*}{\begin{minipage}[b]{\linewidth}\raggedright
\textbf{Evolved Factor Pool}
\end{minipage}} &
\multirow{15}{*}{\begin{minipage}[b]{\linewidth}\raggedright
29.63\%
\end{minipage}} \\
\begin{minipage}[b]{\linewidth}\raggedright
\textbf{ARR: 7.06\% \textbar{} RankICIR: 0.166 \textbar{}}
\end{minipage} & \begin{minipage}[b]{\linewidth}\raggedright
\begin{quote}
\textbf{ARR: 7.35\% \textbar{} RankICIR: 0.170 \textbar{} MDD: 9.67\%}
\end{quote}
\end{minipage} \\
\begin{minipage}[b]{\linewidth}\raggedright
\textbf{MDD: 10.7\%}
\end{minipage} & \begin{minipage}[b]{\linewidth}\raggedright
\begin{quote}
\textbf{{[}Hypothesis{]}}
\end{quote}
\end{minipage} \\
\begin{minipage}[b]{\linewidth}\raggedright
\textbf{{[}Hypothesis{]}}
\end{minipage} & \begin{minipage}[b]{\linewidth}\raggedright
\begin{quote}
\textbf{Volatility regime stress amplifies aligned nested momentum
signals.}
\end{quote}
\end{minipage} \\
\begin{minipage}[b]{\linewidth}\raggedright
\textbf{Aligned 1D/5D/20D momentum}
\end{minipage} &
\multirow{2}{*}{\begin{minipage}[b]{\linewidth}\raggedright
\begin{quote}
\textbf{{[}Key Factor{]} VolRegime\_Momentum\_Composite\_5\_1\_5\_20}
\end{quote}
\end{minipage}} \\
\begin{minipage}[b]{\linewidth}\raggedright
\textbf{predicts sustained price}
\end{minipage} \\
\multirow{2}{*}{\begin{minipage}[b]{\linewidth}\raggedright
\textbf{continuation.}
\end{minipage}} & \begin{minipage}[b]{\linewidth}\raggedright
\begin{quote}
\textbf{Expr:} ZSCORE(TS\_STD(\$return, 5) * (SIGN(TS\_MEAN(\$return,
\end{quote}
\end{minipage} \\
& \begin{minipage}[b]{\linewidth}\raggedright
\begin{quote}
1)) + SIGN(TS\_MEAN(\$return, 5)) + SIGN(TS\_MEAN(\$return,
\end{quote}
\end{minipage} \\
\begin{minipage}[b]{\linewidth}\raggedright
\textbf{{[}Key Factor{]}}
\end{minipage} & \begin{minipage}[b]{\linewidth}\raggedright
\begin{quote}
20))))
\end{quote}
\end{minipage} \\
\begin{minipage}[b]{\linewidth}\raggedright
\textbf{Nested\_Momentum\_Alignment\_}
\end{minipage} & \begin{minipage}[b]{\linewidth}\raggedright
\begin{quote}
\textbf{{[}Feedback{]}}
\end{quote}
\end{minipage} \\
\begin{minipage}[b]{\linewidth}\raggedright
\textbf{1\_5\_20}
\end{minipage} & \begin{minipage}[b]{\linewidth}\raggedright
\begin{quote}
\textbf{Over-engineered \& Lacks Generalization. Drastically simplify.}
\end{quote}
\end{minipage} \\
\begin{minipage}[b]{\linewidth}\raggedright
\textbf{Expr:} SIGN(TS\_MEAN(\$return,
\end{minipage} & \begin{minipage}[b]{\linewidth}\raggedright
\begin{quote}
\textbf{Replace complex nested logic with a standard 5D/20D Vol Ratio +}
\end{quote}
\end{minipage} \\
\begin{minipage}[b]{\linewidth}\raggedright
1)) * SIGN(TS\_MEAN(\$return,
\end{minipage} & \begin{minipage}[b]{\linewidth}\raggedright
\begin{quote}
\textbf{20D Momentum. Use linear weighting instead of complex
conditions}
\end{quote}
\end{minipage} \\
\begin{minipage}[b]{\linewidth}\raggedright
5)) * SIGN(TS\_MEAN(\$return,
\end{minipage} &
\multirow{2}{*}{\begin{minipage}[b]{\linewidth}\raggedright
\begin{quote}
\textbf{to reduce overfitting.}
\end{quote}
\end{minipage}} \\
\begin{minipage}[b]{\linewidth}\raggedright
20)) + 0.5 *
\end{minipage} \\
\midrule()
\endhead
\bottomrule()
\end{longtable}

\begin{quote}
(TS\_CORR(\$return, \$return,\\
5) + TS\_CORR(\$return,\\
\$return, 20))
\end{quote}

\begin{longtable}[]{@{}
  >{\raggedright\arraybackslash}p{(\columnwidth - 6\tabcolsep) * \real{0.2500}}
  >{\raggedright\arraybackslash}p{(\columnwidth - 6\tabcolsep) * \real{0.2500}}
  >{\raggedright\arraybackslash}p{(\columnwidth - 6\tabcolsep) * \real{0.2500}}
  >{\raggedright\arraybackslash}p{(\columnwidth - 6\tabcolsep) * \real{0.2500}}@{}}
\toprule()
\begin{minipage}[b]{\linewidth}\raggedright
\begin{quote}
\textbf{{[}Feedback{]}}
\end{quote}
\end{minipage} &
\multirow{20}{*}{\begin{minipage}[b]{\linewidth}\raggedright
19.14\%
\end{minipage}} &
\multirow{2}{*}{\begin{minipage}[b]{\linewidth}\raggedright
22.38\%
\end{minipage}} &
\multirow{2}{*}{\begin{minipage}[b]{\linewidth}\raggedright
\begin{quote}
\textbf{Iter 5} \textbar{} \textbf{Crossover}
\end{quote}
\end{minipage}} \\
\begin{minipage}[b]{\linewidth}\raggedright
\begin{quote}
\textbf{Missing Magnitude \& Risk}
\end{quote}
\end{minipage} \\
\begin{minipage}[b]{\linewidth}\raggedright
\begin{quote}
\textbf{Control. Need to upgrade from}
\end{quote}
\end{minipage} & &
\multirow{5}{*}{\begin{minipage}[b]{\linewidth}\raggedright
\begin{quote}
\textbf{Iter 1} \textbar{} \textbf{Initialization}
\end{quote}
\end{minipage}} & \begin{minipage}[b]{\linewidth}\raggedright
\begin{quote}
\textbf{ARR: 7.80\% \textbar{} RankICIR: 0.193 \textbar{} MDD: 11.4\%}
\end{quote}
\end{minipage} \\
\begin{minipage}[b]{\linewidth}\raggedright
\begin{quote}
\textbf{simple sign alignment to}
\end{quote}
\end{minipage} & & & \begin{minipage}[b]{\linewidth}\raggedright
\begin{quote}
\textbf{{[}Hypothesis{]}}
\end{quote}
\end{minipage} \\
\begin{minipage}[b]{\linewidth}\raggedright
\begin{quote}
\textbf{Magnitude + Direction}
\end{quote}
\end{minipage} & & & \begin{minipage}[b]{\linewidth}\raggedright
\begin{quote}
\textbf{Volatility-weighted fusion of Institutional Momentum}
\end{quote}
\end{minipage} \\
\begin{minipage}[b]{\linewidth}\raggedright
\textbf{consistency, gated by volatility to}
\end{minipage} & & &
\multirow{2}{*}{\begin{minipage}[b]{\linewidth}\raggedright
\begin{quote}
\textbf{and Retail Herding Reversals.}
\end{quote}
\end{minipage}} \\
\multirow{2}{*}{\begin{minipage}[b]{\linewidth}\raggedright
\begin{quote}
\textbf{prevent false signals during}
\end{quote}
\end{minipage}} \\
& & \multirow{2}{*}{\begin{minipage}[b]{\linewidth}\raggedright
\begin{quote}
\textbf{ARR: 5.22\% \textbar{} RankICIR: 0.158 \textbar{} MDD: 7.67\%}
\end{quote}
\end{minipage}} &
\multirow{2}{*}{\begin{minipage}[b]{\linewidth}\raggedright
\textbf{{[}Key Factor{]} Institutional\_Momentum\_Score\_20D}
\end{minipage}} \\
\begin{minipage}[b]{\linewidth}\raggedright
\begin{quote}
\textbf{stress.}
\end{quote}
\end{minipage} \\
\multirow{11}{*}{\begin{minipage}[b]{\linewidth}\raggedright
13.27\%
\end{minipage}} & & \begin{minipage}[b]{\linewidth}\raggedright
\begin{quote}
\textbf{{[}Initial Direction{]}}
\end{quote}
\end{minipage} & \begin{minipage}[b]{\linewidth}\raggedright
\begin{quote}
\textbf{Expr:} RANK(TS\_CORR(DELTA(\$close, 1)/\$close,
\end{quote}
\end{minipage} \\
& & \begin{minipage}[b]{\linewidth}\raggedright
\begin{quote}
\textbf{Does STD5/WVMA5 Filter enhance ROC60 \& KLOW?}
\end{quote}
\end{minipage} & \begin{minipage}[b]{\linewidth}\raggedright
\begin{quote}
DELTA(\$volume, 1)/\$volume, 20) *
\end{quote}
\end{minipage} \\
& & \multirow{2}{*}{\begin{minipage}[b]{\linewidth}\raggedright
\begin{quote}
\textbf{{[}Hypothesis{]}}
\end{quote}
\end{minipage}} & \begin{minipage}[b]{\linewidth}\raggedright
\begin{quote}
TS\_MEAN((\$close - \$open)/\$close, 5))
\end{quote}
\end{minipage} \\
& & & \multirow{2}{*}{\begin{minipage}[b]{\linewidth}\raggedright
\begin{quote}
\textbf{{[}Observation{]}}
\end{quote}
\end{minipage}} \\
& & \begin{minipage}[b]{\linewidth}\raggedright
\begin{quote}
\textbf{High Price Vol / Low Volume Vol = Noise Regime. Filters for}
\end{quote}
\end{minipage} \\
& & \multirow{2}{*}{\begin{minipage}[b]{\linewidth}\raggedright
\begin{quote}
\textbf{enhanced Mean Reversion efficacy.}
\end{quote}
\end{minipage}} & \begin{minipage}[b]{\linewidth}\raggedright
\begin{quote}
\textbf{Stronger Signal but Higher Risk. Returns (7.8\%) and}
\end{quote}
\end{minipage} \\
& & & \multirow{2}{*}{\begin{minipage}[b]{\linewidth}\raggedright
\begin{quote}
\textbf{IC improved significantly, but Max Drawdown}
\end{quote}
\end{minipage}} \\
& & \multirow{2}{*}{\begin{minipage}[b]{\linewidth}\raggedright
\begin{quote}
\textbf{{[}Key Factor{]} RegimeFiltered\_Reversal\_5D}
\end{quote}
\end{minipage}} \\
& & & \begin{minipage}[b]{\linewidth}\raggedright
\begin{quote}
\textbf{worsened (-11.4\%). The factors capture the trend but}
\end{quote}
\end{minipage} \\
& & \begin{minipage}[b]{\linewidth}\raggedright
\begin{quote}
\textbf{Expr:} (TS\_STD(\$return, 5) / (TS\_STD(\$return *
\end{quote}
\end{minipage} & \begin{minipage}[b]{\linewidth}\raggedright
\textbf{lack Regime-Adaptive Weighting to protect against}
\end{minipage} \\
& & \begin{minipage}[b]{\linewidth}\raggedright
\begin{quote}
\$volume, 5) + 1e-8)) * TS\_PCTCHANGE(\$close, 60)
\end{quote}
\end{minipage} & \begin{minipage}[b]{\linewidth}\raggedright
\begin{quote}
\textbf{downside volatility.}
\end{quote}
\end{minipage} \\
\midrule()
\endhead
\bottomrule()
\end{longtable}

Figure 6: A case study of iterative hypothesis and factor updates in
QuantaAlpha.

23

\end{document}
